% ===== 00_frontmatter.tex =====
\documentclass[11pt,a4paper]{article}
\usepackage[T1]{fontenc}
\usepackage[utf8]{inputenc}
\usepackage{lmodern}
\usepackage[a4paper,margin=25mm,headheight=14pt]{geometry}
\usepackage{amsmath,amssymb,amsthm,mathtools}
\usepackage{booktabs,longtable,tabularx,array}
\usepackage{microtype}
\usepackage{xcolor}
\usepackage{enumitem}
\usepackage{xurl}
\usepackage{fancyhdr}
\usepackage{hyperref}
\definecolor{inkblue}{RGB}{24,57,82}
\hypersetup{colorlinks=true,linkcolor=inkblue,citecolor=inkblue,urlcolor=inkblue,
  pdftitle={Higher Harmonic Reconstruction and Exact Collision Identities},
  pdfauthor={Research continuation prepared for the ProveIt project},
  pdfsubject={Polylogarithms, Gamma functions, harmonic tails, Stieltjes constants, exact proofs}}
\newtheorem{theorem}{Theorem}[section]
\newtheorem{lemma}[theorem]{Lemma}
\newtheorem{proposition}[theorem]{Proposition}
\newtheorem{corollary}[theorem]{Corollary}
\theoremstyle{definition}
\newtheorem{definition}[theorem]{Definition}
\newtheorem{conjecture}[theorem]{Conjecture}
\theoremstyle{remark}
\newtheorem{remark}[theorem]{Remark}
\numberwithin{equation}{section}
\setlist{itemsep=4pt,topsep=6pt}
\setlength{\emergencystretch}{2em}
\setlength{\parskip}{3pt plus 1pt}
\allowdisplaybreaks[2]
\pagestyle{fancy}
\fancyhf{}
\fancyhead[L]{\small ProveIt research continuation}
\fancyhead[R]{\small Exact higher identities}
\fancyfoot[C]{\thepage}
\renewcommand{\headrulewidth}{0.3pt}
\newcommand{\file}[1]{\nolinkurl{#1}}
\newcommand{\CT}{\operatorname{CT}}
\newcommand{\FP}{\operatorname{FP}}
\title{\textbf{Higher Harmonic Reconstruction\\and Exact Collision Identities}\\[5pt]
  \large Gap ten, three-cluster polylogarithms,\\and Gamma--Stieltjes primitives}
\author{Research continuation prepared for\\Vladimir Reshetnikov's \emph{ProveIt} project}
\date{12 October 2026}
\begin{document}
\maketitle
\begin{abstract}
We continue several explicitly open directions in the inspected ProveIt
manuscript and incoming research collection. First, we evaluate the local
correction for collisions of arbitrarily many polygamma factors along
arbitrary positive paths. For nonresonant monomial scales this becomes a
finite zeta polynomial. With one trigamma factor and any number of centered
digammas, the complete constant is given by a single coefficient formula;
the triple constants are $\zeta(3),0,-\zeta(3)$ in scale order. An explicit
resonant integer-scale triple has a different, fully evaluated answer.
Second, a mixed derivative of a logarithmic simplex face with any number
of singly labelled displaced masses reduces to a divided difference of
one classical polylogarithm. Two integrations give a finite depth-three
colored-polylogarithm dictionary for every three-cluster face $B^r,C,D$,
and hence an ordinary convergent spectral formula for its Gamma analogue.
Third, we prove the complete gap-ten harmonic-tail reconstruction family.
An analytic argument covers every $b\ge73728$, and an exact finite-field
certificate covers all $73722$ remaining indices. A general Green-function
formula represents inverse coefficients by determinants of bounded size
and yields ordinary convergent polygamma-product sums. Finally, classical
Euler homotopy converts the inherited universal normal coefficients into
nonrecursive Gamma and all-index Stieltjes antiderivative certificates.
We prove local analytic convergence of both infinite-jet series and their
translation identity. Proof dependencies, computational certificates,
scope corrections and further research questions are recorded explicitly.
\end{abstract}
\noindent\textbf{Keywords.} Multiple polylogarithms; Hurwitz zeta;
generalized Stieltjes constants; polygamma collisions; harmonic-tail
determinants; Euler homotopy; finite parts; exact computation.

\clearpage
{\setlength{\parskip}{0pt}\small\tableofcontents}
\clearpage

% ===== 00_overview.tex =====
\section{Purpose, principal results and conventions}
\label{intro:section}

The aim is to add proved identities to the current research programme,
while preserving the distinctions between an ordinary integral, a specified
finite part, a formal reconstruction theorem and a numerical observation.
The starting points are the manuscript and incoming directories named by
the project maintainer \cite{RepoManuscript,RepoIncoming}. All twelve ZIP
archives present in the retrieved incoming listing were downloaded and
their mathematical source files inspected by topic. Ninety-two manuscript
chapter files were also retrieved. This is a targeted proof and status
audit of the directions used here, not a claim to have checked every line
of every chapter. File hashes and the retrieval scope are supplied in
\file{data/source_manifest.json}.

The newest source matters. In particular, the report
\emph{Polygamma Moments and Stieltjes Resonances} already completes gap
eight, treats all positive monomial scales for two Stieltjes factors,
and evaluates a substantial joint-polygamma range
\cite{PolygammaResonances}. The expanded incoming version of
\emph{Universal Gamma Normal Forms} already gives the coefficient of
every canonical monomial \cite{UniversalIncoming}. Those results are
inputs here. The new work starts with the next fixed gap, additional
collision factors, a third mass cluster, and the still missing
nonrecursive primitive and infinite-jet convergence statements.

\subsection{A map of the contributions}

\begin{longtable}{@{}>{\raggedright\arraybackslash}p{0.21\textwidth}>{\raggedright\arraybackslash}p{0.38\textwidth}>{\raggedright\arraybackslash}p{0.34\textwidth}@{}}
\toprule
Direction & Result proved here & Precise remaining limit\\
\midrule
\endhead
Polygamma collisions & A finite Hermite-divided-difference correction
for any number of factors and arbitrary positive shifts;
Theorem~\ref{pc:uniform}. & The general merged product is retained as a
specified finite part. Positive Stieltjes indices require further work.\\[5pt]
Evaluated constants & The complete one-trigamma, centered-digamma family,
Theorem~\ref{pc:all-powers}; nonresonant and resonant triples. & The scale
hypotheses determine which constant is selected.\\[5pt]
Multiple mass clusters & A classical-polylogarithm divided difference
for every labelled mixed derivative, and a depth-three primitive for
$B^r,C,D$; Theorems~\ref{tri:thm:mixed} and~\ref{tri:thm:three}.
& The Gamma result is an ordinary spectral sum; no general finite
arithmetic closure is proved.\\[5pt]
Harmonic reconstruction & Every unequal even pair with gap ten has the
optimal formal prefix; Theorem~\ref{gt:main}. & The unrestricted
unequal-even-pair conjecture remains open.\\[5pt]
Inverse coefficients & An $(m+1)$-square bordered determinant at fixed
half-gap $m$; Theorem~\ref{gt:inverse}. & Its entries still use a finite
causal recurrence.\\[5pt]
Gamma--Stieltjes primitives & A nonrecursive finite formula for every
compensator $H_m$ and locally convergent normal and primitive series;
Theorems~\ref{hp:closed}, \ref{hp:analyticN} and~\ref{hp:analyticH}.
& Sharp convergence domains, coefficient signs and arithmetic
denominators remain separate questions.\\
\bottomrule
\end{longtable}

The contact and Gamma-primitive arguments are analytic or algebraic
proofs at every order. The gap-ten theorem has a finite exact
computational component, fully exposed below. Numerical checks of
polylogarithms and shifted integrals are diagnostics of the displayed
formulas; they are not substitutes for the proofs.

\subsection{Two representative identities}

The most compact collision identity can be stated before introducing
the full local correction. Let $\alpha_1>\alpha_2>\alpha_3>0$ be
linearly independent over $\mathbb Q$, and let $r\in\{1,2,3\}$.
Then
\begin{equation}\label{intro:triple}
\boxed{\displaystyle
\CT_{\epsilon\downarrow0}\int_0^1
 \psi'(x+\epsilon^{\alpha_r})
 \prod_{j\ne r}\psi(x+\epsilon^{\alpha_j})\,dx
 =\begin{cases}
  \zeta(3),&r=1,\\
  0,&r=2,\\
  -\zeta(3),&r=3.
 \end{cases}}
\end{equation}
Here the smallest shift carries the largest exponent. Every integral
is ordinary for $\epsilon>0$; the operation outside it is a specified
asymptotic constant. The fully evaluated subtraction formula in
Section~\ref{pc:section} also expresses these constants as ordinary
renormalized limits. No fitting of divergent coefficients is needed.

For the higher family put $\Psi=\psi+\gamma$ and
$f(x)=x(\psi(1+x)+\gamma)=\sum_{k\ge2}(-1)^k\zeta(k)x^k$.
With the analogous nonresonant ordering of $N$ scales, the complete
constant with the derivative in position $r$ is
\begin{equation}\label{intro:allN}
\boxed{\displaystyle
T_{N,r}=-[x^N]\left\{\frac{(f-1)^N}{N}
                     +f^{N-r}(f-1)^{r-1}\right\}.}
\end{equation}
Thus a single coefficient formula yields a weight-$N$ zeta polynomial
for every $N$ and every derivative position. Its adjacent differences
and its zero row sum are identities, rather than patterns inferred from
the first rows.

The representative spectral identity has a different character. For
the explicitly defined ordinary simplex face $I_k$ and Gamma field $G_a$
of Section~\ref{tri:sec:main}, all $r,m\ge1$ satisfy
\begin{equation}\label{intro:mixed}
\partial_{C_1}\cdots\partial_{C_m}
 I_{r+m}(G_a;B^r,C_1,\ldots,C_m)
 =(-1)^r r!\,\mathfrak g_{r,a,B}[C_1,\ldots,C_m],
\end{equation}
where
\[
\mathfrak g_{r,a,B}(C)=\sum_{n\ge0}
 \frac{\operatorname{Li}_r((C-B)/((n+a)^2+C))}{(n+a)^2+C}.
\]
The number of independent masses does not increase the depth of the
polylogarithm in the differentiated formula. The finite three-cluster
primitive reconstructs the information on the coordinate edges as well
as the common mixed derivative.

\subsection{Conventions used throughout}

We write $\psi^{(r)}$ for the $r$th argument derivative of the digamma
function, with $\psi^{(0)}=\psi$. The Hurwitz function is
$\zeta(s,a)$, and a superscript on $\zeta$ denotes an order derivative
unless an argument derivative is explicitly displayed. The Stieltjes
normalization is
\begin{equation}\label{intro:Stieltjes}
\zeta(1+s,a)=\frac1s+
    \sum_{q\ge0}\frac{(-1)^q}{q!}\gamma_q(a)s^q,
\qquad \gamma_0(a)=-\psi(a).
\end{equation}
These conventions agree with the standard Gamma and Hurwitz-zeta
identities \cite{DLMF,DLMFHurwitz}.
Generalized harmonic numbers are
$H_n^{(s)}=\sum_{j=1}^n j^{-s}$, with $H_0^{(s)}=0$.
The Bernoulli convention is $B_1=-1/2$.

For a locally finite power--logarithm expansion,
$\CT_\epsilon$ selects the coefficient of
$\epsilon^0(\log\epsilon)^0$. The merged finite part at the lower
endpoint is always taken in the displayed coordinate $x$:
$\CT_{\delta\downarrow0}\int_\delta^1$.
It is not invariant under an arbitrary change of cutoff coordinate.
All endpoint subtractions and all shifts are therefore kept explicit.

For divided differences, one node gives evaluation and repeated nodes
give derivatives divided by the corresponding factorial. Cumulative
colors for multiple polylogarithms are defined in
\eqref{tri:eq:cumulative}; this convention makes the relevant real
branches and zero-color limits transparent. Mass derivatives are taken
in independently labelled variables before imposing collisions.

Finally, formal rank in the harmonic reconstruction problem means
rational matrix rank on declared symbols. It does not imply arithmetic
independence of the numerical zeta or moment values assigned to them.
The new conclusions are established relative to the inspected project
sources. Classical homotopy, divided-difference and finite-part methods
are credited; a bounded literature search cannot certify worldwide
priority for the assembled identities.

% ===== 01_collisions.tex =====
% Mathematical fragment for integration by the parent agent.
% Standard theorem environments, amsmath/amssymb, and \psi are assumed.
% Every label begins with pc:.

\section{Evaluated contacts for arbitrarily many polygamma factors}
\label{pc:section}

The latest incoming \emph{Polygamma Moments and Stieltjes Resonances}
report, its further-research question~6, asks for three independently shifted
Stieltjes factors \cite{PolygammaResonances}.  The following result answers the local-contact problem in the entire
Stieltjes-index-zero sector, with any number of factors and arbitrary
nonnegative argument derivative orders.  Its uniform comparison permits
all positive paths, including nonmonomial paths.  Spectral derivatives
corresponding to positive Stieltjes indices are not claimed here.
The finite-part integration and divided-difference tools are classical
\cite{DeBoor,TicaGalapon,VillanuevaGalapon};
the claimed contribution relative to the inspected sources is their
complete singular-subset assembly and the evaluated consequences below.

\subsection{The rational kernel as one Hermite divided difference}

For positive nodes $a_1,\ldots,a_p$, positive integers
$h_1,\ldots,h_p$, and $H=h_1+\cdots+h_p$, define, for
$0\le d\le H-1$,
\begin{equation}\label{pc:kernel-def}
 K_{\boldsymbol h,d}(\boldsymbol a)
 =\lim_{R\to\infty}\left\{
       \int_0^R\frac{x^d\,dx}{\prod_{i=1}^p(x+a_i)^{h_i}}
                 -\mathbf1_{d=H-1}\log R\right\}.
\end{equation}
Thus only the last numerator degree needs a finite part at infinity.
For a smooth function $F$ on the positive axis, write
$[a_1^{(h_1)},\ldots,a_p^{(h_p)}]F$ for its Hermite divided
difference with the node $a_i$ repeated $h_i$ times; its order is $H-1$.
The standard confluent convention \cite{DeBoor} gives
$[a^{(H)}]F=F^{(H-1)}(a)/(H-1)!$.

\begin{lemma}[Completely evaluated local kernel]\label{pc:kernel}
For every indicated degree and every positive set of nodes,
\begin{equation}\label{pc:kernel-dd}
 \boxed{\displaystyle
 K_{\boldsymbol h,d}(\boldsymbol a)
     =(-1)^{H+d}
       [a_1^{(h_1)},\ldots,a_p^{(h_p)}](z^d\log z).}
\end{equation}
In particular this is a finite expression in rational functions of the
nodes and their logarithms, including all confluent cases.
\end{lemma}

\begin{proof}
First take $H$ distinct nodes $b_1,\ldots,b_H$ and all multiplicities
equal to one.  The partial fraction coefficients of
$x^d/\prod_j(x+b_j)$ are
\[
 c_i=\frac{(-b_i)^d}{\prod_{j\ne i}(b_j-b_i)}.
\]
If $d\le H-2$ their sum is zero; if $d=H-1$ their sum is one.
Integrating from zero to $R$ and subtracting the prescribed logarithm
therefore gives
\[
 -\sum_i c_i\log b_i
   =(-1)^{H+d}\sum_i
          \frac{b_i^d\log b_i}{\prod_{j\ne i}(b_i-b_j)},
\]
which is exactly the divided difference.  Coalescing nodes proves the
general case.  For $d<H-1$ dominated convergence applies directly.
For $d=H-1$, apply it to the integral on $(0,1)$ and to
$x^d/\prod_j(x+b_j)-1/x$ on $(1,\infty)$, whose tail has a
locally uniform $O(x^{-2})$ bound.  This also proves the confluent
continuation without choosing singular individual coefficients.
\end{proof}

For direct implementation at distinct nodes one can instead use
\begin{align}\label{pc:partial-fractions}
 c_{i,q}&=[y^{h_i-q}]\,(y-a_i)^d
                      \prod_{j\ne i}(y+a_j-a_i)^{-h_j},\\
 K_{\boldsymbol h,d}(\boldsymbol a)
   &=-\sum_i c_{i,1}\log a_i
       +\sum_i\sum_{q=2}^{h_i}
                    \frac{c_{i,q}a_i^{1-q}}{q-1}.
                    \label{pc:kernel-rational}
\end{align}
Every coefficient in the first line is finite.  These equations are
also a transparent way to extract power--logarithm constants along
specified paths.

\subsection{The complete uniform local correction}

Fix $N\ge1$ and $r_1,\ldots,r_N\in\mathbb Z_{\ge0}$.  Put
\[
 h_i=r_i+1,\qquad p_i=(-1)^{r_i+1}r_i!,\qquad
 g_i(x)=\psi^{(r_i)}(1+x).
\]
The classical polygamma recurrence \cite[Eq.~5.15.5]{DLMF} gives
\begin{equation}\label{pc:split}
 \psi^{(r_i)}(x)=p_i x^{-h_i}+g_i(x).
\end{equation}
For a nonempty subset $S\subseteq\{1,\ldots,N\}$ let
\begin{align}
 H_S&=\sum_{i\in S}h_i,\qquad p_S=\prod_{i\in S}p_i,\\
 A_{S,d}(\boldsymbol a)
   &=[x^d]\prod_{j\notin S}g_j(x+a_j)
     =\sum_{\substack{d_j\ge0\ (j\notin S)\\\sum d_j=d}}
          \prod_{j\notin S}
              \frac{\psi^{(r_j+d_j)}(1+a_j)}{d_j!}.
              \label{pc:A}
\end{align}
The empty product convention makes $A_{S,0}=1$ and $A_{S,d}=0$
for $d>0$ when $S$ is the full set.  Crucially, the shifts in
\eqref{pc:A} are retained exactly.

Define the fixed-coordinate merged finite part
\begin{equation}\label{pc:merged}
 F_{\boldsymbol r}
    =\operatorname{CT}_{\delta\downarrow0}
                  \int_\delta^1\prod_{i=1}^N\psi^{(r_i)}(x)\,dx
\end{equation}
and the completely evaluated correction
\begin{align}\label{pc:Q}
 Q_{\boldsymbol r}(\boldsymbol a)
   &=\sum_{\varnothing\ne S\subseteq\{1,\ldots,N\}}
        p_S\sum_{d=0}^{H_S-1}A_{S,d}(\boldsymbol a)
                        K_{\boldsymbol h_S,d}(\boldsymbol a_S)\notag\\
   &=\sum_{\varnothing\ne S}
       \left(\prod_{i\in S}r_i!\right)
       \sum_{d=0}^{H_S-1}(-1)^d A_{S,d}(\boldsymbol a)
        [a_i^{(h_i)}:i\in S](z^d\log z).
\end{align}

\begin{theorem}[Every polygamma collision with positive shifts]
\label{pc:uniform}
As all $a_i\downarrow0$, with no restriction on their relative rates,
\begin{equation}\label{pc:uniform-formula}
 \boxed{\displaystyle
 \int_0^1\prod_{i=1}^N\psi^{(r_i)}(x+a_i)\,dx
          =F_{\boldsymbol r}+Q_{\boldsymbol r}(\boldsymbol a)+o(1).}
\end{equation}
The $o(1)$ is uniform as $\max_i a_i\to0$ through positive shifts.
The same statement holds after replacing any zero-order factor
$\psi$ by $\psi+c_i$, provided the same constant is added to its
regular part in \eqref{pc:A} and in the merged finite part.
\end{theorem}

\begin{proof}
Expand the product by choosing its singular part on a subset $S$.
For $S\ne\varnothing$, put
\[
 R_S(x;\boldsymbol a)=\prod_{i\in S}(x+a_i)^{-h_i},\qquad
 G_S(x;\boldsymbol a)=\prod_{j\notin S}g_j(x+a_j).
\]
Uniform Taylor expansion on a fixed compact interval gives
\[
 G_S(x;\boldsymbol a)
    =\sum_{d=0}^{H_S-1}A_{S,d}(\boldsymbol a)x^d
                   +x^{H_S}B_S(x;\boldsymbol a),
\]
where $B_S$ is uniformly bounded and tends to $B_S(x;0)$ for each
$x\in[0,1]$.  Since $R_S(x;\boldsymbol a)\le x^{-H_S}$ for
positive shifts, the integral of $x^{H_S}R_SB_S$ tends uniformly
to $\int_0^1B_S(x;0)dx$.  This estimate does not compare one shift
with another.

For $d\le H_S-2$, comparison with the tail on $(1,\infty)$ yields
\[
 \int_0^1x^dR_S(x;\boldsymbol a)dx
       =K_{\boldsymbol h_S,d}(\boldsymbol a_S)
                              -\frac1{H_S-d-1}+o(1).
\]
For $d=H_S-1$, compare instead with
$\int_1^\infty(x^{H_S-1}R_S-1/x)dx$; this tends to zero uniformly,
so the same formula holds with its rational term omitted.  Both
tail comparisons follow from an $O(\max a_i)x^{-2}$ bound after
the appropriate leading power is removed.

The empty-subset term converges ordinarily.  The limits of all
remainders and all displayed rational terms sum exactly to
$F_{\boldsymbol r}$: a merged monomial $x^{d-H_S}$ contributes
$-1/(H_S-d-1)$ for $d\le H_S-2$, while $x^{-1}$ contributes
zero to the fixed-coordinate constant on $(0,1)$.  This proves
\eqref{pc:uniform-formula}.  Uniform dominated convergence follows
by taking the supremum over the shrinking parameter cube in the
same integrable bound.  Constants added to the regular parts
do not change the proof.
\end{proof}

This theorem evaluates the entire local correction.  It does not assert
a new scalar evaluation of every merged product $F_{\boldsymbol r}$;
some such merged values remain spectral or harmonic constants in the
prior manuscript.  The next family eliminates that baseline by an
ordinary antiderivative and hence gives complete scalar identities.

\subsection{Nonresonant scales and their ordering}

Assume
\begin{equation}\label{pc:scales}
 a_i=\lambda_i\epsilon^{\alpha_i},\qquad
 \alpha_1>\alpha_2>\cdots>\alpha_N>0,\qquad\lambda_i>0,
\end{equation}
and that $\alpha_1,\ldots,\alpha_N$ are linearly independent over
$\mathbb Q$.  For a nonempty subset $S$, write $m(S)=\max S$;
its node is the largest shift within $S$.  Let $H_0=0$ and
$H_r=1+\cdots+1/r$.

\begin{theorem}[Evaluated nonresonant constant]\label{pc:nonresonant}
The ordinary shifted integral has a power--logarithm expansion through
its constant, and
\begin{equation}\label{pc:nonresonant-constant}
 \boxed{\displaystyle
 \operatorname{CT}_{\epsilon}
      \int_0^1\prod_i\psi^{(r_i)}(x+a_i)dx
  =F_{\boldsymbol r}
    -\sum_{\varnothing\ne S}p_S A_{S,H_S-1}(0)
            \bigl(\log\lambda_{m(S)}+H_{r_{m(S)}}\bigr).}
\end{equation}
The formula includes the same centering modification as
Theorem~\ref{pc:uniform}.  Its correction is a finite polynomial in
ordinary zeta values, Euler's constant and the displayed logarithms.
\end{theorem}

\begin{proof}
The rational formula \eqref{pc:kernel-rational} expands in convergent
geometric series in ratios of smaller shifts to larger shifts.
After each finite pole order is factored out, all remaining increments
of the exponent are positive differences $\alpha_i-\alpha_j$ or
positive regular-part exponents.  Therefore only finitely many terms
can occur below any fixed power.  Analytic expansion of the regular
coefficients has the same property.

For $d<H_S-1$, the rational coefficients of $K_{\boldsymbol h_S,d}$
are homogeneous of degree $d+1-H_S<0$, and the logarithms are
$\log\lambda_i+\alpha_i\log\epsilon$.  A zero exponent would,
by rational independence, require every integer exponent of every
shift variable to be zero.  Powers coming from regular-part shifts
occur only in variables outside $S$, so they must all be zero.
The remaining degree cannot be zero, a contradiction.  Thus no
noncritical degree contributes a constant.

For $d=H_S-1$, regular-part powers outside $S$ again cannot contribute.
Put $b=a_{m(S)}$ and $\rho_i=a_i/b$.  Scaling the defining integral
gives
\[
 K_{\boldsymbol h_S,H_S-1}(\boldsymbol a_S)
     =-\log b+
          K_{\boldsymbol h_S,H_S-1}
                         ((\rho_i)_{i\ne m(S)},1).
\]
Dominated convergence on $(0,1)$ and on the compensated tail gives
\begin{align*}
 \lim_{\rho_i\to0}
 K_{\boldsymbol h_S,H_S-1}((\rho_i),1)
 &=\lim_{R\to\infty}\left\{
   \int_0^R\frac{x^{h_{m(S)}-1}}{(1+x)^{h_{m(S)}}}dx-\log R
                                        \right\}\\
 &=-H_{h_{m(S)}-1}.
\end{align*}
For the first dominated-convergence step use the majorant
$x^{h_{m(S)}-1}/(1+x)^{h_{m(S)}}$ on $(0,1)$; the compensated
tail is uniformly $O(x^{-2})$.  The last evaluation follows from
$t=x/(1+x)$ and
$t^{h-1}/(1-t)=1/(1-t)-\sum_{j=0}^{h-2}t^j$.
Consequently the critical constant is
$-\log\lambda_{m(S)}-H_{r_{m(S)}}$.  Substitution into
\eqref{pc:Q} proves the theorem.
\end{proof}

Rational independence is a sufficient clean hypothesis, not an assertion
that every rationally dependent choice is resonant.  The exact kernel
formula remains valid for every positive path and permits finite
coefficient extraction in every specified monomial case.

\subsection{A family with complete zeta-valued answers}

Let
\[
 \Psi(x)=\psi(x)+\gamma,\qquad
 h(x)=\psi(1+x)+\gamma,\qquad
 f(x)=xh(x)=\sum_{k\ge2}(-1)^k\zeta(k)x^k.
\]
The series is the classical digamma Taylor expansion
\cite[Eq.~5.7.4]{DLMF}.
Thus $\Psi(x)=-1/x+h(x)$ and $\Psi(1)=h(0)=0$.
Take unit amplitudes in \eqref{pc:scales}.  For $1\le r\le N$,
define the full constant of the ordinary interval integral
\begin{equation}\label{pc:T-def}
 T_{N,r}=\operatorname{CT}_{\epsilon}
       \int_0^1\psi'(x+\epsilon^{\alpha_r})
               \prod_{j\ne r}\Psi(x+\epsilon^{\alpha_j})\,dx.
\end{equation}

\begin{theorem}[Every centered one-derivative moment]
\label{pc:all-powers}
For every $N\ge2$ and $1\le r\le N$,
\begin{equation}\label{pc:all-powers-formula}
 \boxed{\displaystyle
 T_{N,r}=-[x^N]\left\{
           \frac{(f(x)-1)^N}{N}
               +f(x)^{N-r}(f(x)-1)^{r-1}\right\}.}
\end{equation}
This is a finite polynomial in ordinary zeta values of total weight $N$.
It is independent of the particular nonresonant exponents and depends
only on the position of the derivative in their ordering.  Moreover
\begin{equation}\label{pc:sum-jump}
 \sum_{r=1}^N T_{N,r}=0,\qquad
 T_{N,r+1}-T_{N,r}
       =[x^N]f(x)^{N-r-1}(f(x)-1)^{r-1}.
\end{equation}
\end{theorem}

\begin{proof}
The merged antiderivative is $\Psi(x)^N/N$, so its finite part is
\[
 B_N=-\frac1N[x^N](f(x)-1)^N.
\]
In \eqref{pc:nonresonant-constant}, only a singular subset that
contains the derivative, has that derivative as its largest shift,
and contains any selection of its $r-1$ smaller shifts can contribute:
all other subsets have harmonic number $H_0=0$.
A subset of size $s$ has $H_S=s+1$, prefactor $p_S=(-1)^{s-1}$,
and regular coefficient $[x^s]h(x)^{N-s}$.  Its contact is therefore
$(-1)^s[x^s]h(x)^{N-s}$, with multiplicity $\binom{r-1}{s-1}$.
Their sum is
\begin{align*}
 C_{N,r}
 &=\sum_{s=1}^r(-1)^s\binom{r-1}{s-1}
                           [x^s]h(x)^{N-s}\\
 &=-[x^N]f(x)^{N-r}(f(x)-1)^{r-1}.
\end{align*}
Adding $B_N$ gives the theorem.  The adjacent difference is immediate.
For the sum use
$\sum_{r=1}^N f^{N-r}(f-1)^{r-1}=f^N-(f-1)^N$ and
$[x^N]f^N=0$, since $f$ starts in degree two.
\end{proof}

The first rows, simplified with $\zeta(2)^2=\tfrac52\zeta(4)$, are
\begin{align}
 (T_{2,1},T_{2,2})&=(0,0),\\
 (T_{3,1},T_{3,2},T_{3,3})&=(\zeta(3),0,-\zeta(3)),
                                                  \label{pc:triple-row}\\
 (T_{4,1},T_{4,2},T_{4,3},T_{4,4})
   &=\frac{\zeta(4)}4(-11,-1,5,7),
                                                  \label{pc:fourth-row}\\
 (T_{5,1},\ldots,T_{5,5})
   &=\bigl(\zeta(5)-4\zeta(2)\zeta(3),
            \zeta(5)-4\zeta(2)\zeta(3),\notag\\
   &\qquad\zeta(5)-2\zeta(2)\zeta(3),
            2\zeta(2)\zeta(3),
            8\zeta(2)\zeta(3)-3\zeta(5)\bigr).
                                                  \label{pc:fifth-row}
\end{align}
These are full asymptotic constants of the integrals, not merely their
differences from the merged finite part.

\subsection{The uncentered triple and an explicit resonant path}

For the uncentered triple the merged value is
\begin{equation}\label{pc:triple-baseline}
 \operatorname{CT}_{\delta\downarrow0}
       \int_\delta^1\psi'(x)\psi(x)^2\,dx
           =\zeta(3)-2\gamma\zeta(2).
\end{equation}
The three nonresonant contacts, with unit amplitudes, are respectively
$2\gamma\zeta(2)$,
$2\gamma\zeta(2)-\zeta(3)$, and
$2\gamma\zeta(2)-2\zeta(3)$.  Thus the full constants remain
exactly the triple row \eqref{pc:triple-row}.  A concrete admissible
choice is $(\alpha_1,\alpha_2,\alpha_3)=(\sqrt3,\sqrt2,1)$.
The cyclic sum is zero, also directly forced by integrating the
derivative of $\prod_i\psi(x+a_i)$ and taking its constant.

For comparison, the resonant path $(a,b,c)=(\epsilon^3,\epsilon^2,
\epsilon)$ with the derivative at $a$ has the different exact value
\begin{equation}\label{pc:integer-triple}
 \boxed{\displaystyle
 \operatorname{CT}_{\epsilon}
   \int_0^1\psi'(x+\epsilon^3)\psi(x+\epsilon^2)
                            \psi(x+\epsilon)dx
 =4+2\gamma-\zeta(2)+\zeta(3)
       +\left(\frac92-\gamma\right)\zeta(4)
       +\zeta(5)-\zeta(6).}
\end{equation}
Here is a finite certificate of its constant extraction.
At unit amplitudes every logarithm in \eqref{pc:kernel-rational}
is an integer multiple of $\log\epsilon$, so only its rational part
can contribute to the log-free constant.  Any singular subset not
containing the derivative has only simple poles and hence contributes
zero.  The four other subsets give the following constants, where
$h(x)=\psi(1+x)$:
\begin{center}
\begin{tabular}{c|l}
Singular subset & Its contact constant\\\hline
$\{a\}$ & $-\gamma\zeta(4)+\zeta(2)^2+2\gamma\zeta(2)$\\
$\{a,b\}$ & $\gamma-\zeta(3)+2\zeta(4)+\zeta(5)-\zeta(6)$\\
$\{a,c\}$ & $\gamma-\zeta(2)+\zeta(3)$\\
$\{a,b,c\}$ & $4$
\end{tabular}
\end{center}
For example the last rational part is
$1/[a(b-a)(c-a)]=\epsilon^{-6}/[(1-\epsilon)(1-\epsilon^2)]$,
whose constant is four.  The two-factor subset with remaining
regular factor $h(x+c)$ has rational part
\[
 -\frac{h(c)-a h'(c)+a^2h''(c)/2}{a(b-a)},
\]
and the exchanged expression gives the other row.  The singleton
has rational part
$h(b)h(c)/a-h'(b)h(c)-h(b)h'(c)$.
Taking the indicated finite Taylor coefficients and adding
\eqref{pc:triple-baseline} proves \eqref{pc:integer-triple}.

\paragraph{Ordinary renormalized limits.}
The symbol $\operatorname{CT}_{\epsilon}$ always selects
$\epsilon^0(\log\epsilon)^0$ from a proved expansion; it is not the
ordinary limit of a divergent integral.  All the results also give
ordinary limits: subtract the finite expression
$Q_{\boldsymbol r}(\boldsymbol a)$ from the original integral and
then add its evaluated constant.  Equation~\eqref{pc:uniform-formula}
proves the resulting limit equals the stated full constant.  Thus the
subtraction is an explicit finite formula in polygammas, rational
functions and logarithms, with no fitted coefficients or new integrals.


The exact and numerical verification records for this section are
described in Section~\ref{repro:section}. The kernel methods are
classical \cite{DeBoor,Lesch,TicaGalapon,VillanuevaGalapon};
the complete assembly and evaluated families are the continuation
results relative to the inspected project sources.

% ===== 02_three_clusters.tex =====
% Scratch research section. All labels carry the prefix tri:.
% Requires amsmath, amssymb, amsthm; no custom macros assumed.
% The face I_k is the ordinary perspective face used in the incoming reports.

\section{Mixed mass derivatives and the three-cluster polylogarithm dictionary}
\label{tri:sec:main}

The incoming two-cluster theorem evaluates
$I_{r+s}(G_a;B^r,C^s)$, and explicitly leaves $B^r,C,D$ with two
unequal displaced masses as its next target \cite{GammaColored}. We resolve that target
by a finite dictionary of weight $r+2$ and depth at most three,
uniformly in $r$. An accompanying theorem is stronger at the
differential level: the mixed derivative in arbitrarily many singly
labelled displaced masses is a divided difference of one classical
polylogarithm. The spectral sum in the Gamma application remains an
ordinary infinite sum; no arithmetic closure or minimal-depth claim
is made.

\subsection{Normalization and the inherited coefficient identity}

For $k\ge2$, write $\Sigma_k=\{t_i\ge0:\sum_i t_i=1\}$ with
measure $dt_1\cdots dt_{k-1}$. For nonempty $T\subseteq\{1,\ldots,k\}$
let $t_T=\sum_{i\in T}t_i$ and
$b_T=\sum_{i\in T}t_i b_i/t_T$. The perspective $t_TF(b_T)$ is
filled by zero when $t_T=0$. Set
\begin{align}
 \mathcal H_k(F;t)&=\sum_{\varnothing\ne T}
                (-1)^{k-|T|}t_TF(b_T),\label{tri:eq:H}\\
 I_k(F;b_1,\ldots,b_k)&=\int_{\Sigma_k}
                   \mathcal H_k(F;t)\frac{dt}{t_1\cdots t_k}.
                   \label{tri:eq:I}
\end{align}
These integrals are ordinary and absolutely convergent for $C^1$
functions on the mass interval. The perspectives are Lipschitz and
the alternating numerator vanishes on each coordinate face, so it
is bounded by a constant times $\min_i t_i$. This also supplies
locally integrable bounds after sufficiently small independent
complex powers of the $t_i$ are inserted. Constants and affine
functions are annihilated.

Let $D_F(\boldsymbol\alpha;\boldsymbol b)$ be the normalized
Dirichlet mean of $F(\sum t_i b_i)$, with total parameter
$A=\sum_i\alpha_i$. The independent-regulator identity inherited
from the two-cluster report \cite{GammaColored} implies
\begin{equation}\label{tri:eq:coefficient}
 I_k(F;\boldsymbol b)
  =[\alpha_1\cdots\alpha_k]\{A D_F(\boldsymbol\alpha;\boldsymbol b)\}.
\end{equation}
Indeed, the normalized alternating Dirichlet integral equals
$(\prod_i\alpha_i)^{-1}$ times the subset alternating sum of
$A_TD_F(\boldsymbol\alpha_T;\boldsymbol b_T)$; every proper subset
omits a variable and disappears under squarefree coefficient
extraction. When $r$ labels have the common mass $B$, their
parameters aggregate into $X$, and hence
\begin{equation}\label{tri:eq:cluster-coefficient}
 I_{r+m}(F;B^r,C_1,\ldots,C_m)
  =r![X^rY_1\cdots Y_m]\{A D_F(X,Y_1,\ldots,Y_m)\}.
\end{equation}
Here $r\ge1$, $m\ge1$, and $A=X+Y_1+\cdots+Y_m$.
All coefficients below are taken in the holomorphic continuation
of the displayed normalized mean. For the logarithms used below,
its existence also follows directly by expanding around the largest
positive spectral factor: the resulting Dirichlet moment series
has a geometric convergence margin and denominators
$(1+A)_{n-1}$.

\subsection{Arbitrarily many mixed derivatives reduce to one classical polylogarithm}

For an analytic function $f$, its divided difference \cite{DeBoor} is denoted by
$f[C_1,\ldots,C_m]$; thus one node means $f(C_1)$ and two nodes
mean $(f(C_1)-f(C_2))/(C_1-C_2)$. Repeated nodes are assigned
their unique confluent values, so
$f[C,\ldots,C]=f^{(m-1)}(C)/(m-1)!$.

\begin{theorem}[Classical-polylogarithm mixed derivatives]
\label{tri:thm:mixed}
Let $\sigma>0$, $B,C_1,\ldots,C_m\ge0$, $r,m\ge1$, and put
\begin{equation}\label{tri:eq:g}
 g_{r,\sigma,B}(C)=\frac{1}{\sigma+C}
       \operatorname{Li}_r\!\left(\frac{C-B}{\sigma+C}\right).
\end{equation}
Then
\begin{equation}\label{tri:eq:mixed}
 \boxed{\displaystyle
 \partial_{C_1}\cdots\partial_{C_m}
 I_{r+m}\bigl(U\mapsto-\log(\sigma+U);
                  B^r,C_1,\ldots,C_m\bigr)
 =(-1)^r r!\,g_{r,\sigma,B}[C_1,\ldots,C_m].}
\end{equation}
The derivative is taken in each independently labelled displaced
mass before any nodes are made equal. In particular, the bound of
one on classical polylogarithm depth is independent of both $r$ and
$m$.
\end{theorem}

\begin{proof}
Write $S_B=\sigma+B$ and $S_j=\sigma+C_j$. Differentiate the
Dirichlet mean in \eqref{tri:eq:cluster-coefficient}. Since
\[
 \partial_{C_1}\cdots\partial_{C_m}
   [-\log(S_Bt_0+\sum_jS_jt_j)]
  =(-1)^m(m-1)!\frac{t_1\cdots t_m}
                      {(S_Bt_0+\sum_jS_jt_j)^m},
\]
the coefficient of $Y_1\cdots Y_m$, first at $\Re X>0$, is
\[
 (-1)^m(m-1)!X\int_{\Sigma_{m+1}}
        \frac{t_0^{X-1}}{(S_Bt_0+\sum_jS_jt_j)^m}\,dt.
\]
This follows either by shifting each $Y_j$ by one in the Dirichlet
density or directly from $1/\Gamma(Y_j)=Y_j+O(Y_j^2)$.
At a fixed $t=t_0$, the ordinary simplex resolvent integral gives
\[
 \int_{\Sigma_m}\frac{du}
       {(S_Bt+(1-t)\sum_jS_ju_j)^m}
  =\frac{1}{(m-1)!\prod_{j=1}^m(S_Bt+(1-t)S_j)}.
\]
For completeness, this is the positive-parameter Feynman identity:
write the product of $m$ reciprocal positive numbers as a Laplace
integral over the positive orthant and separate its radial variable.
Therefore the required coefficient of the differentiated mean is
\begin{equation}\label{tri:eq:parameter-integral}
 (-1)^mX\int_0^1
  \frac{t^{X-1}(1-t)^{m-1}}
       {\prod_j(S_Bt+(1-t)S_j)}\,dt.
\end{equation}
At fixed $t$, the divided difference of
$C\mapsto(S_Bt+(1-t)(\sigma+C))^{-1}$ is
$(-1)^{m-1}(1-t)^{m-1}$ divided by the product in
\eqref{tri:eq:parameter-integral}. Thus that expression is the
negative divided difference of
\[
 R_X(C)=X\int_0^1
             \frac{t^{X-1}}{S_Bt+(1-t)(\sigma+C)}\,dt.
\]
Writing $z=(C-B)/(\sigma+C)$ yields, initially for $|z|<1$,
\begin{equation}\label{tri:eq:RX}
 R_X(C)=\frac1{\sigma+C}
       \left(1+\sum_{n=1}^{\infty}\frac{Xz^n}{X+n}\right),
 \qquad
 [X^r]R_X(C)=\frac{(-1)^{r-1}}{\sigma+C}\operatorname{Li}_r(z).
\end{equation}
The series is normally convergent near $X=0$ on each compact
$z$-disk. Its integral representation continues it through every
real $z<1$ without crossing a singularity; equivalently, subtract
the value at $t=0$ before continuing the $X$ integral. Now take
$r![X^r]$ and use \eqref{tri:eq:cluster-coefficient}. All mass
derivatives and confluent limits are justified by local holomorphy
away from nonpositive spectral factors. This proves the theorem.
\end{proof}

For $m=1$, integration from $C=B$ recovers the inherited
one-displaced formula
\begin{equation}\label{tri:eq:one-displaced}
 I_{r+1}(-\log(\sigma+U);B^r,C)
   =(-1)^r r!\operatorname{Li}_{1,r}
                         \!\left(\frac{C-B}{\sigma+C}\right),
\end{equation}
because $d\operatorname{Li}_{1,r}(z)/dz=
\operatorname{Li}_r(z)/(1-z)$ and the face vanishes when all masses
coincide. This recovery checks the factorial and sign conventions.

\subsection{A uniform finite dictionary for the three-cluster primitive}

It is useful to specify colored multiple polylogarithms by cumulative
colors, which avoids apparent quotients by zero. For positive
indices $s_1,\ldots,s_d$, define
\begin{equation}\label{tri:eq:cumulative}
 \mathcal C_{s_1,\ldots,s_d}(c_1,\ldots,c_d)
 =\sum_{n_1>\cdots>n_d\ge1}
 \frac{c_1^{n_1-n_2}c_2^{n_2-n_3}\cdots
            c_{d-1}^{n_{d-1}-n_d}c_d^{n_d}}
      {n_1^{s_1}\cdots n_d^{s_d}}.
\end{equation}
For $|c_j|<1$ this series converges absolutely. In the conventional
colored notation it is
$\operatorname{Li}_{s_1,\ldots,s_d}
(c_1,c_2/c_1,\ldots,c_d/c_{d-1})$, interpreted continuously
when a color is zero. The one-variable convention used by the
incoming reports is
$\operatorname{Li}_{s_1,\ldots,s_d}(z)=
\mathcal C_{s_1,\ldots,s_d}(z,\ldots,z)$.
For real cumulative colors below one, the continuation is the
iterated-integral branch continued from zero, with letters
$0,1/c_1,\ldots,1/c_d$ and the path $[0,1]$. No letter crosses
that path on the stated real domain.

Put $H(z)=\operatorname{Li}_1(z)=-\log(1-z)$ and
\begin{align}
 \mathcal S_r(z,w)
   ={}&\mathcal C_{1,1,r}(z,w,w)
      +\mathcal C_{1,r,1}(z,z,w)
      +\mathcal C_{1,r+1}(z,w),\label{tri:eq:S}\\
 \mathcal K_r(z,w)
   ={}&H(w)\operatorname{Li}_{1,r}(z)
       +H(z)\operatorname{Li}_{1,r}(w)
       -\mathcal S_r(z,w)-\mathcal S_r(w,z),\label{tri:eq:K}\\
 \mathcal P_r(z,w)
   ={}&(r+1)\{\operatorname{Li}_{1,r+1}(z)
                    +\operatorname{Li}_{1,r+1}(w)\}
       +\mathcal K_r(z,w).\label{tri:eq:P}
\end{align}
Every term has weight $r+2$; the products in \eqref{tri:eq:K}
can be shuffled as iterated integrals into depth at most three,
retaining the same safe letters for real negative colors.
Thus the formula gives
a finite dictionary with an upper bound of three on depth for
every $r\ge1$.

\begin{theorem}[Three-cluster logarithmic face]
\label{tri:thm:three}
For $\sigma>0$, $B,C,D\ge0$, $r\ge1$, let
$z=(C-B)/(\sigma+C)$ and $w=(D-B)/(\sigma+D)$. Then
\begin{equation}\label{tri:eq:three}
 \boxed{\displaystyle
 I_{r+2}(-\log(\sigma+U);B^r,C,D)
       =(-1)^{r+1}r!\mathcal P_r(z,w).}
\end{equation}
In particular, if $B\le C,D$, all cumulative colors in the finite
dictionary lie in $[0,1)$, regardless of the sizes of the mass
displacements. The result therefore requires no small-displacement
hypothesis.
\end{theorem}

\begin{proof}
The key harmonic identity, valid for $|z|,|w|<1$, is
\begin{equation}\label{tri:eq:K-series}
 \mathcal K_r(z,w)
  =\sum_{m,n\ge1}\frac{z^mw^n}{mn}
    \{H_{m-1}^{(r)}+H_{n-1}^{(r)}-H_{m+n-1}^{(r)}\}.
\end{equation}
To check its finite reduction, the first two harmonic terms give
the two products in \eqref{tri:eq:K}. For the last term use
$1/(mn)=(1/m+1/n)/(m+n)$ and put $N=m+n$. The harmonic index
$j<N$ can lie below $m$, above $m$, or equal $m$; these three
cases give respectively
$\mathcal C_{1,1,r}(z,w,w)$,
$\mathcal C_{1,r,1}(z,z,w)$, and
$\mathcal C_{1,r+1}(z,w)$. Exchanging $z,w$ supplies the other
half. All rearrangements are absolute.

Write $f_r(z)=(1-z)\operatorname{Li}_r(z)$. Differentiating
\eqref{tri:eq:K-series} gives
\begin{equation}\label{tri:eq:K-mixed}
 \partial_z\partial_w\mathcal K_r(z,w)
   =-\frac{f_r[z,w]}{(1-z)(1-w)}.
\end{equation}
Indeed,
$\sum_{n\ge1}H_{n-1}^{(r)}z^{n-1}=
\operatorname{Li}_r(z)/(1-z)$, while the diagonal sum involving
$H_{m+n-1}^{(r)}$ equals
\[
 \frac{\operatorname{Li}_r(z)/(1-z)
                    -\operatorname{Li}_r(w)/(1-w)}{z-w}.
\]
Combining these expressions proves \eqref{tri:eq:K-mixed},
including its removable diagonal.

Now $z'=(1-z)/(\sigma+C)$,
$w'=(1-w)/(\sigma+D)$,
\[
 z-w=\frac{(\sigma+B)(C-D)}{(\sigma+C)(\sigma+D)},
 \qquad f_r(z)=(\sigma+B)g_{r,\sigma,B}(C).
\]
Consequently the mixed $C,D$ derivative of the right side of
\eqref{tri:eq:three} agrees exactly with
Theorem~\ref{tri:thm:mixed} for $m=2$. On $C=B$ or $D=B$,
$\mathcal K_r$ vanishes and the remaining term is precisely
the inherited one-displaced formula with $r+1$ base labels.
Equality follows from two integrations of the common mixed
derivative, starting on those two coordinate edges. The proof
initially applies in a neighborhood of $C=D=B$ and throughout
$B\le C,D$ by the absolutely convergent color series. Analytic
continuation along the positive mass domain gives every stated
$B,C,D\ge0$.
\end{proof}

The formula is symmetric under $C\leftrightarrow D$. Its two
principal collision checks are
\begin{align}
 \mathcal P_r(z,0)&=(r+1)\operatorname{Li}_{1,r+1}(z),
                      \label{tri:eq:axis}\\
 \mathcal P_r(z,z)&=2r\operatorname{Li}_{1,r+1}(z)
       +2\sum_{j=1}^{r}\operatorname{Li}_{1,j,r+1-j}(z).
                      \label{tri:eq:diagonal}
\end{align}
The second identity recovers twice the predecessor's $P_{r,2}$.
It also follows directly from the shuffle identity
\[
 H(z)\operatorname{Li}_{1,r}(z)
  =\operatorname{Li}_{1,1,r}(z)
      +\operatorname{Li}_{1,r,1}(z)
      +\sum_{j=1}^{r}\operatorname{Li}_{1,j,r+1-j}(z).
\]
Thus the three-cluster formula has the correct labelled
multiplicity even when the two displaced masses become equal.

\subsection{Gamma faces, confluent derivatives, and polygamma values}

For $a>0$, $U\ge0$, put
\[
 G_a(U)=\log\Gamma(a+i\sqrt U)+\log\Gamma(a-i\sqrt U)
                                     -\log(2\pi).
\]
The analytic logarithms are continued from positive real Gamma
arguments. The ordinary Gamma product \cite[\S5.8]{DLMF} gives
\[
 G_a(U)=2\log\Gamma(a)-\log(2\pi)
       -\sum_{n=0}^{\infty}
           \log\left(1+\frac{U}{(n+a)^2}\right).
\]
Its second mass derivative is normally convergent on compact
nonnegative mass sets, so every ordinary face may be applied
termwise after subtracting an affine function from each term.

\begin{corollary}[The three-cluster Gamma dictionary]
\label{tri:cor:Gamma}
For $a>0$, $B,C,D\ge0$ and $r\ge1$, set
\[
 z_n=\frac{C-B}{(n+a)^2+C},\qquad
 w_n=\frac{D-B}{(n+a)^2+D}.
\]
Then
\begin{equation}\label{tri:eq:Gamma}
 \boxed{\displaystyle
 I_{r+2}(G_a;B^r,C,D)
   =(-1)^{r+1}r!\sum_{n=0}^{\infty}\mathcal P_r(z_n,w_n).}
\end{equation}
The sum is ordinary and absolutely convergent, locally normally
in the stated parameter domain. Each spectral term has a finite
colored multiple-polylogarithm representation of depth at most
three.
\end{corollary}

\begin{proof}
Apply Theorem~\ref{tri:thm:three} termwise to the Gamma product.
For a direct tail bound, \eqref{tri:eq:K-series} gives
\[
 \mathcal P_r(z,w)=\frac{r+1}{2}(z^2+w^2)-zw
                     +O((|z|+|w|)^3).
\]
Both colors are $O(n^{-2})$, so the spectral summands are
$O(n^{-4})$. Compact local holomorphic extensions give the same
normal convergence for every fixed parameter derivative.
\end{proof}

Define the normally convergent classical-polylogarithm sum
\begin{equation}\label{tri:eq:spectralg}
 \mathfrak g_{r,a,B}(C)=\sum_{n=0}^{\infty}
      \frac{\operatorname{Li}_r((C-B)/((n+a)^2+C))}
           {(n+a)^2+C}.
\end{equation}
Theorem~\ref{tri:thm:mixed} gives, for every $m\ge1$,
\begin{equation}\label{tri:eq:Gamma-mixed}
 \boxed{\displaystyle
 \partial_{C_1}\cdots\partial_{C_m}
 I_{r+m}(G_a;B^r,C_1,\ldots,C_m)
   =(-1)^r r!\mathfrak g_{r,a,B}[C_1,\ldots,C_m].}
\end{equation}
For example, the first case beyond the two-cluster primitive is
\begin{equation}\label{tri:eq:dilog-example}
 \partial_C\partial_D I_4(G_a;B,B,C,D)
   =\frac{2}{C-D}\sum_{n=0}^{\infty}
      \left\{\frac{\operatorname{Li}_2(z_n)}{(n+a)^2+C}
             -\frac{\operatorname{Li}_2(w_n)}{(n+a)^2+D}\right\}.
\end{equation}
At $C=D$ it is the corresponding derivative of
$\mathfrak g_{2,a,B}$; no singular mass denominator remains.

More generally, making $m$ independent displaced masses equal
after differentiation gives
\begin{equation}\label{tri:eq:confluent}
 \left.\partial_{C_1}\cdots\partial_{C_m}
 I_{r+m}(G_a;B^r,C_1,\ldots,C_m)\right|_{C_j=C}
 =\frac{(-1)^r r!}{(m-1)!}\,
                    \partial_C^{m-1}\mathfrak g_{r,a,B}(C).
\end{equation}
Every finite derivative here involves only classical polylogarithms
$\operatorname{Li}_{r-j}$ and rational mass factors. Negative
polylogarithm orders are rational functions. The distinction
between these labelled mixed derivatives and total derivatives
along a diagonal is essential.

At the complete mass collision $C=B$, the formula has a finite
harmonic-polygamma specialization. Let
$H_n^\star(\{1\}^r)=\sum_{n\ge j_1\ge\cdots\ge j_r\ge1}
(j_1\cdots j_r)^{-1}$. For $m\ge2$,
\begin{equation}\label{tri:eq:complete-collision}
 \left.\partial_{C_1}\cdots\partial_{C_m}
 I_{r+m}(G_a;B^r,C_1,\ldots,C_m)\right|_{C_j=B}
 =(-1)^{r+m}r!H_{m-1}^\star(\{1\}^r)
             \sum_{n=0}^{\infty}((n+a)^2+B)^{-m}.
\end{equation}
Indeed the coefficient of $u^N$ in
$(1+u)^{-1}\operatorname{Li}_r(u/(1+u))$ is
$(-1)^{N+1}H_N^\star(\{1\}^r)$; this follows from the elementary
binomial identity
$\sum_{j=1}^N(-1)^{j-1}\binom Nj/j^r=H_N^\star(\{1\}^r)$.
At $B=0$ the final sum is $\zeta(2m,a)$; at $B>0$ it is a
finite derivative of the paired digamma resolvent
\[
 \sum_{n=0}^{\infty}((n+a)^2+B)^{-1}
  =\frac{\psi(a+i\sqrt B)-\psi(a-i\sqrt B)}{2i\sqrt B}.
\]
This is a consistency check with the predecessor's all-degree
collision formula, not a separate novelty claim for those
coefficients.

\subsection{An independent ordinary integral for verification}

The finite dictionary also admits a single scalar logarithmic
integral that is convenient for numerical verification. Put
$q=(B-C)/(\sigma+B)$, $p=(B-D)/(\sigma+B)$,
$L(t)=\log(1-t)$, $a(t)=-\log(1-qt)$, and
$b(t)=-\log(1-pt)$. Then
\begin{align}\label{tri:eq:log-integral}
 I_{r+2}(-\log(\sigma+U);B^r,C,D)
  =\int_0^1\bigg[&L^r\{a'(L+b)+b'(L+a)\}\notag\\
       &\quad +(r+1)\frac{L^r(a+b)}t
       +r\frac{L^{r-1}ab}t\bigg]dt.
\end{align}
All apparent singularities at zero are removable or integrable;
at one only powers of $\log(1-t)$ occur. One derivation uses
independent Gamma variables and Frullani's formula to write
\[
 A D_{-\log}(X,Y,Z)
   =\mathcal L(X,Y,Z)
     +\int_0^1(1-t)^{X+Y+Z}
       \frac{d}{dt}\frac{(1-qt)^{-Y}(1-pt)^{-Z}-1}{t}\,dt.
\]
Here $\mathcal L$ is linear in $X,Y,Z$.
Extract $r![X^rYZ]$ and integrate the double-pole term using
the primitive $(1-t)/t$, which vanishes at $t=1$.
This produces \eqref{tri:eq:log-integral} without endpoint
finite parts.

For an additional check wholly in the local mass series, let
$h_j(N)=H_{N-1}^{\star}(\{1\}^j)$. When $|q|,|p|<1$,
the independent Dirichlet moment expansion gives
\begin{align}\label{tri:eq:collision-check}
 I_{r+2}(-\log(\sigma+U);B^r,C,D)
 =(-1)^r r!\bigg[&\sum_{m,n\ge1}\frac{q^mp^n}{mn}h_r(m+n)\notag\\
 &-(r+1)\sum_{m\ge1}\frac{q^m+p^m}{m}h_{r+1}(m)\bigg].
\end{align}
The two sums arise respectively from terms using both displaced
Dirichlet parameters and from terms using just one. This provides
an independently organized harmonic computation of the same face.

The supplied verifier records 231 assertions: 192 exact rational
binomial--harmonic checks, fourteen comparisons from seven scalar
parameter cases, 24 independently computed mixed-derivative checks,
and one Gamma spectral check. The scalar cases include
$r=1,2,3,6$, a displaced mass below the base, an equal-outlier
collision, and a zero displacement. At $\sigma=1$, $B=0$,
$C=1/5$, $D=2/5$, $r=2$, for example, the common scalar value is
\[
 -0.307903203351030388542590020916578537605405701313403\ldots.
\]
The finite colored dictionary, the ordinary scalar integral, and
the local Dirichlet series agree at 75-digit working precision
with scaled residuals below $1.5\cdot10^{-75}$. The independent
mixed-derivative checks use the predecessor's monomial-face
formula and include both complete and partial confluence; their
largest scaled residual is below $6\cdot10^{-72}$.
The Gamma example is
\begin{equation}\label{tri:eq:Gamma-value}
 I_4(G_1;0,0,1/5,2/5)
  =-0.3411810552611179668555780745233632208321622014483386\ldots.
\end{equation}
A finite colored-MPL spectral head plus an exact-form Hurwitz
tail agrees with the independent complete Hurwitz collision
series to a scaled residual below $1.4\cdot10^{-76}$.
These are floating-point diagnostics, not interval certificates
or substitutes for the analytic proofs.

% ===== 03_gap_ten.tex =====
% Parent-ready proof fragment. All labels use gt:.
% Bibliography keys used: OrderedBoundary, PolygammaResonances, BilinearTails, DLMF.
% Current source paths and verification files are listed in README_gap_ten.txt.
\section{The complete gap-ten reconstruction theorem}
\label{gt:section}

The latest incoming report \emph{Polygamma Moments and Stieltjes
Resonances} proves the complete gap-eight family and gives a computable
eventual threshold at every fixed gap \cite{PolygammaResonances}.
Thus the next unresolved fixed-gap family in the inspected material is
the gap of ten. We settle that entire family. The analytic improvement
is to keep different error budgets for odd and even Gamma modes; only
two of the five modes require the larger error. The finite proof uses
exact fast polynomial arithmetic, so a substantially longer interval
can be checked without quadratic propagation of the fundamental array.

We also give a bounded-size formula for every inverse coefficient at
an arbitrary fixed gap. At gap ten it uses determinants of order six,
independently of the two tail orders. This supplies a constructive
answer to one version of the earlier compact-inverse-row question
\cite{BilinearTails}. It retains a finite causal recurrence and does
not assert a short product evaluation of those coefficients.
These advances are relative to the inspected project sources; no
claim of historical priority is made.

\subsection{The original harmonic-tail coordinates}

Let $b\ge6$, $p=2b-10$, $q=2b$, and
\[
 A_v(n)=\zeta(v,n+1)=\zeta(v)-H_n^{(v)},\qquad x_n=n+\tfrac12.
\]
The series
\[
 D_{p,q}(s)=\sum_{n\ge0} A_p(n)A_q(n)x_n^{-s}
\]
converges absolutely for $\operatorname{Re}s>3-p-q$ and has a
meromorphic continuation. Write $D_r=D_{p,q}(-2r)$.
All these values are regular: the centered Hurwitz expansion
\begin{equation}\label{gt:tail-expansion}
 A_v(n)\sim\sum_{h\ge0}c_{v,h}x_n^{1-v-2h},\qquad
 c_{v,0}=\frac1{v-1},\quad
 c_{v,h}=\frac{B_{2h}(1/2)}{(2h)!}(v)_{2h-1}\quad(h\ge1)
\end{equation}
allows subtraction of successive terms, and the resulting possible
poles of $D_{p,q}$ are the odd integers $3-p-q-2h$.
The Bernoulli convention is $B_1=-1/2$, and $(v)_k$ is a rising factorial.

The inherited block transform is
\begin{align}
 C_j&=\sum_{r=0}^{j}\binom{2j+1}{2r}2^{-2(j-r)}D_r,
       \label{gt:C}\\
 D_r&=\sum_{j=0}^{r}\frac{\binom{2r}{2j}}{2j+1}
                      B_{2r-2j}(1/2)C_j.\label{gt:D}
\end{align}
For $j\ge b$ it has the finite late-coordinate form
\begin{equation}\label{gt:late}
 C_j=\sum_{k=0}^{b-1}a_{jk}Z_k,\qquad
 Z_k=\zeta(2k),\qquad Z_0=-\tfrac12.
\end{equation}
For $d=2j+1$ and $e=d-p-q+2k+1$, put $a_{jk}=0$ if $e<0$;
otherwise the exact coefficient is
\begin{equation}\label{gt:source-entry}
 a_{jk}=B_e\left\{
 \frac{\binom{d-q+1}{p-2k}}{d-q+1}
 +\frac{\binom{d-p+1}{q-2k}}{d-p+1}\right\}.
\end{equation}
Each summand is omitted unless its lower binomial index is between
$2$ and its upper index. In particular, a zero lower index is omitted.
This is the exact convention in the original source, and it will be
retained throughout the proof.

For context, the finite summation underlying these formulas follows
from
\[
 A_p(n-1)A_q(n-1)-A_p(n)A_q(n)
 =n^{-p}A_q(n)+n^{-q}A_p(n)+n^{-p-q}.
\]
Multiply by $n^{2j+1}$ and sum by parts. The constant term at
$N\to\infty$ of
\begin{align*}
 &N^{2j+1}A_p(N)A_q(N)
 +\sum_{n=1}^{N} n^{2j+1-q}A_p(n)\\
 &\hspace{20mm}+\sum_{n=1}^{N}n^{2j+1-p}A_q(n)
 +\sum_{n=1}^{N}n^{2j+1-p-q}
\end{align*}
is $C_j$. Expanding the centered difference
$(x+1/2)^{2j+1}-(x-1/2)^{2j+1}$ gives \eqref{gt:C};
the inverse generating factor is $z/(2\sinh(z/2))$, giving
\eqref{gt:D}. Finite Faulhaber summation gives
\eqref{gt:source-entry}. These coordinate identities and their
constant-term interpretation are inherited from the ordered-boundary
development \cite{OrderedBoundary}.

Define the original $b$-square matrix
\[
 M_{b,5}=(a_{b+r,k})_{0\le r,k<b}.
\]
On the specified formal labels
\begin{equation}\label{gt:formal-labels}
 C_0,\ldots,C_{b-1},Z_0,\ldots,Z_{b-1},
\end{equation}
the first $2b$ block rows have coefficient matrix
$\operatorname{diag}(I_b,M_{b,5})$.
The diagonal of \eqref{gt:D} is $1/(2r+1)$, so the corresponding
coefficient determinant of $D_0,\ldots,D_{2b-1}$ is
\begin{equation}\label{gt:prefix-det}
                       \frac{\det M_{b,5}}{(4b-1)!!}.
\end{equation}
Formal rank below always means rational matrix rank on
\eqref{gt:formal-labels}. It does not assert independence after
these labels are evaluated as numerical constants.

\begin{theorem}[Every gap-ten pair has the optimal prefix]
\label{gt:main}
For every integer $b\ge6$, the matrix $M_{b,5}$ is nonsingular.
Consequently, for every even $p\ge2$ and $q=p+10$, the first $q$
centered moments
\[
 D_{p,q}(0),D_{p,q}(-2),\ldots,D_{p,q}(-2(q-1))
\]
have formal rank $q$ and generate all later centered even moments
over $\mathbb Q$. The number $q$ is optimal for the $q$ specified
formal coordinates. Their coefficient determinant is
\eqref{gt:prefix-det}.
\end{theorem}

\subsection{A rational recurrence of fixed dimension}

Put
\begin{align*}
 \beta_n&=\frac{2|B_{2n}|}{(2n)!}\quad(n\ge1),\\
 a_u&=0\quad(1\le u\le5),\qquad
 a_u=(2u)^{\underline{10}}\quad(u>5),\\
 b_s&=(2s-1)^{\underline{10}},\qquad
 d_s=2(a_s+b_s)\quad(s>5).
\end{align*}
The underline denotes a falling factorial; the symbols $b_s$ here
are recurrence factors, not Bernoulli numbers. Every $d_s$ is positive.
Define the rational array
\begin{align}
 R_{u,j}&=\delta_{uj}&&(1\le u,j\le5),\notag\\
 R_{s,j}&=\frac1{d_s}\sum_{u=1}^{s-1}
             \beta_{s-u}(b_s+a_u)R_{u,j}&&(s>5),
 \label{gt:recurrence}
\end{align}
and let
\[
 F_{b,5}=(R_{b+i,j})_{1\le i,j\le5},\qquad f_{b,5}=\det F_{b,5}.
\]
The inherited complementary-minor identity is
\begin{equation}\label{gt:bridge}
 \det M_{b,5}=
 \frac{(-1)^{b(b-1)/2}}{4^bb!}
            \left(\prod_{s=6}^{b+5}d_s\right)f_{b,5}.
\end{equation}
Its normalization is important. For an even series
$R(z)=\sum_{n\ge1}r_nz^{2n}$, the coefficient of $z^{2s-11}$ in
\begin{equation}\label{gt:null}
 \mathcal T_5R=(R^{(10)}-R^{(10)}(0))\cot(z/2)
                               +(R\cot(z/2))^{(10)}
\end{equation}
is
\[
 d_sr_s-\sum_{u<s}\beta_{s-u}(b_s+a_u)r_u.
\]
The assignment $a_5=0$ records the subtracted derivative, equivalently
the omitted lower binomial index zero in \eqref{gt:source-entry}.
Truncate the five fundamental null series after degree $2b$ and
apply the complementary-minor identity to this coefficient operator.
Reversing the original $b$ columns and restoring their factorial
scales gives exactly \eqref{gt:bridge}; a full normalization proof
is in \cite[Gamma-basis section]{OrderedBoundary}.

\subsection{The inherited Gamma basis, with the normalization made explicit}

Let
\[
 W(x)=\left|\Gamma\!\left(\tfrac14+\tfrac{ix}{2}\right)\right|^{20},
 \qquad
 P_j(x)=\sinh(\pi(5-j)x)\ (1\le j<5),\qquad
 P_5(x)=\tanh(\pi x).
\]
Write
\begin{equation}\label{gt:t}
 t_{n,j}=\frac4{(2n)!}\operatorname{Im}
                  \int_0^\infty(x+i/2)^{2n}P_j(x)W(x)\,dx.
\end{equation}
These are absolutely convergent ordinary integrals. The Gamma-basis
theorem \cite{OrderedBoundary} states that
\[
 G_j(z)=4\sin(z/2)\int_0^\infty\sinh(xz)P_j(x)W(x)\,dx
\]
is a null germ of \eqref{gt:null}, and that
\begin{equation}\label{gt:gamma-bridge}
 (t_{b+i,j})_{i,j=1}^{5}=F_{b,5}T_5,
 \qquad T_5=(t_{n,j})_{n,j=1}^{5},\qquad \det T_5>0.
\end{equation}

Here is the analytic mechanism, to specify precisely what is inherited.
For an allowed multiplier $P$, set
\[
 Q_P(t)=2\sinh(t/2)\int_{\mathbb R}e^{ixt}P(x)W(x)\,dx.
\]
The Fourier contours in the two terms obtained by differentiating
$Q_P$ ten times can be shifted by $\pm i$. Gamma recurrence gives
\[
 (iy-\tfrac12)^{10}W(y+i)
 =(iy+\tfrac12)^{10}W(y-i)
 =2^{10}\left[
    \Gamma(\tfrac34+iy/2)\Gamma(\tfrac34-iy/2)\right]^{10}.
\]
The nearest poles of the common Gamma factor are at $\pm3i/2$,
outside the two shifted contours. Also $P(y+i)=P(y-i)$.
The sinh modes therefore have no contact term. For tanh the two
simple poles at $y=\pm i/2$, each with residue $1/\pi$, give
$Q_P^{(10)}(0)=i2^{12}\pi^5$; the other four top derivatives vanish.
The resulting identity is the hyperbolic version of
\eqref{gt:null}. Rotation $G_j(z)=iQ_{P_j}(iz)$ gives the stated
null equation. Stirling decay justifies all contour shifts and
differentiations: on the vertical boundaries the decay is at least
$e^{-\pi|\operatorname{Re}y|}$ times a fixed polynomial.

For the sign, $\operatorname{Im}(x+i/2)^{2n}$ has leading term
$nx^{2n-1}$. A triangular row operation and continuous Cauchy--Binet
express $\det T_5$ as a positive constant times
\[
 \int_{(0,\infty)^5}
 \Delta(x_1^2,\ldots,x_5^2)
 \Delta(\cosh\pi x_1,\ldots,\cosh\pi x_5)
 \prod_{r=1}^{5}x_rW(x_r)\tanh(\pi x_r)\,d\mathbf x.
\]
Indeed $\sinh(ku)/\tanh u=\cosh u\,U_{k-1}(\cosh u)$,
so the multiplier alternant has column degrees $4,3,2,1,0$;
their reversal has positive sign. Both Vandermondes change sign
together under any interchange, and their product is strictly
positive off the diagonals. This proves $\det T_5>0$.
The coefficient identity in \eqref{gt:gamma-bridge} then follows
from uniqueness in \eqref{gt:recurrence}.

\subsection{An effective phase estimate at half-gap five}

Put
\[
 \lambda_j=j\pi,\qquad A=2^5(2\pi)^{10},\qquad
 a_j^*=A/2\ (j<5),\qquad a_5^*=A.
\]

\begin{lemma}[Specialization of the effective Gamma estimate]
\label{gt:phase}
For every integer $N\ge147456$ and $1\le j\le5$,
\begin{align}
 \left|\frac1{a_j^*}\int_0^\infty(x+i/2)^NP_j(x)W(x)\,dx
 -e^{i\lambda_j/2}\left(
 \frac{\Gamma(N-4)}{\lambda_j^{N-4}}
 +\frac{5i}{2}\frac{\Gamma(N-5)}{\lambda_j^{N-5}}
                         \right)\right|
 <33\frac{\Gamma(N-6)}{\lambda_j^{N-6}}.\label{gt:phase-bound}
\end{align}
\end{lemma}

\begin{proof}
This is the $m=5$ case of the effective phase lemma of
\cite{PolygammaResonances}; a specialized proof follows.
For $x\ge10$, complex Stirling with its remainder \cite[\S5.11(ii)]{DLMF}
gives
\begin{equation}\label{gt:density}
 \frac{P_j(x)W(x)}{a_j^*}
 =x^{-5}e^{-\lambda_jx}(1+\delta_j(x)),\qquad
                      |\delta_j(x)|\le10x^{-2}.
\end{equation}
To check the constants, for $z=1/4+ix/2$ the remainder after
$1/(12z)$ has magnitude at most $1/(90|z|^3)$. With $t=1/(2x)$,
the logarithm of $W/(Ax^{-5}e^{-5\pi x})$ is
\[
 -\frac52\log(1+t^2)
 +5\left(\frac{\arctan t}{t}-1\right)
 +\frac{20}{3}\frac{t^2}{1+t^2}+20\operatorname{Re}\rho(z).
\]
Its absolute value is at most
$65/(24x^2)+16/(9x^3)<15/(4x^2)$.
Using $|e^u-1|\le2|u|$ in this range, and the bound
$2e^{-2\pi x}\le1/(9x^2)$ for the normalized multiplier, proves
\eqref{gt:density}. On $0\le x\le10$ one may use the generous
bound $P_j(x)W(x)/a_j^*<2^{460}$, obtained from
$\Gamma(1/4)<5$, $\pi<4$, and $e<4$.

Write $S_2=\Gamma(N-6)/\lambda_j^{N-6}$.
The density error on $x\ge\sqrt N$ is at most $20S_2$ because
\[
 |x+i/2|^N\le x^N\exp(N/(8x^2))<2x^N.
\]
On $10\le x\le\sqrt N$ its bound is
$\sqrt N(N+1)^{N/2}$.
For the model integral on $[10,\infty)$, put $y=x+i/2$ and
shift the ray down to the positive real axis. No pole is crossed.
The algebraic factor becomes $(1-i/(2y))^{-5}$, and Taylor's
theorem, for $y\ge10$, gives
\[
 \left|(1-i/(2y))^{-5}-1-\frac{5i}{2y}\right|
                         \le\frac{45}{4y^2}.
\]
For example the positive Taylor remainder is at most
$15u^2(1-u)^{-7}<45u^2$ for $u=1/(2y)\le1/20$.
Thus this part of the model error is at most $(45/4)S_2$.

The original compact integral, the short vertical segment, and the
two retained model integrals from $0$ to $10$ are together bounded
by $2^{720}15^N$. To absorb the remaining terms, set $k=N-7$.
Since $\lambda_j<20$ and $k!\ge(k/3)^k$,
\[
 S_2>\frac1{20}\left(\frac{k}{60}\right)^k.
\]
For $N\ge147456$, $k\ge3N/4$ and
$\sqrt{N+1}\le(5/4)\sqrt N$, so
\begin{align*}
 \frac{\sqrt N(N+1)^{N/2}}{S_2}
     &\le N^9(100/\sqrt N)^N<N^9 2^{-N}\le\tfrac14,\\
 \frac{2^{720}15^N}{S_2}
     &\le2^{720}N^8(1200/N)^N
                  <2^{720}N^8 4^{-N}<\tfrac14.
\end{align*}
Here $N^9\le2^{N-2}$ holds at $147456$ and persists by induction.
The total error is less than
$(20+45/4+1/2)S_2<33S_2$, proving \eqref{gt:phase-bound}.
\end{proof}

\subsection{Different determinant budgets for odd and even modes}

\begin{proposition}[A positive tail beginning at an explicit index]
\label{gt:positive-tail}
For every integer $b\ge73728$, one has $f_{b,5}>0$.
\end{proposition}

\begin{proof}
Apply Lemma~\ref{gt:phase} at $N_i=2b+2i$, $1\le i\le5$.
An odd column has leading imaginary phase $\pm1$ and relative
error bounded by
\[
 \frac{33\lambda_j^2}{(N_i-5)(N_i-6)}
                       <\frac{8448}{(N_i-6)^2},
\]
since $5\pi<16$. An even column has zero leading imaginary phase.
Its first nonzero term is the second term in
\eqref{gt:phase-bound}; its relative error is at most
\[
 \frac{66\lambda_j}{5(N_i-6)}
                       <\frac{858}{5(N_i-6)},
\]
because the largest even mode is $4\pi<13$.

Extract the positive row factors
\[
             \frac{4\Gamma(N_i-4)}{\Gamma(N_i+1)}\pi^{-2i}
\]
and positive column factors $a_j^*\lambda_j^{-2b+4}$.
In an even column extract additionally
$5\lambda_j/(2(2b-3))$. The five extracted signs are
$(+,-,-,+,+)$, with positive product. The remaining matrix is
$E_{ij}=j^{-2i}(1+\eta_{ij})$.
For its three odd columns,
\[
 |\eta_{ij}|<\frac{8448}{(2b-4)^2}<2^{-20}.
\]
For the two even columns the additional row factor is
$\rho_i=(2b-3)/(2b+2i-5)$; hence
\[
 |\eta_{ij}|
 <\frac8{2b-3}+\frac{858}{5(2b-4)}
 <\frac{898}{5(2b-4)}<2^{-9}.
\]
Both final inequalities hold already at $b=73728$ and improve
monotonically thereafter.

For $V=(j^{-2i})_{i,j=1}^{5}$, exact rational arithmetic gives
\begin{equation}\label{gt:V}
 \det V=\frac{49}{2764800000000},\qquad
 \operatorname{perm}V=\frac{106852831033}{25798901760000000000},
 \qquad
 \frac{\det V}{\operatorname{perm}V}
       =\frac{457228800}{106852831033}.
\end{equation}
Every permutation product uses exactly two even columns and three
odd columns. Therefore
\begin{align*}
 |\det E-\det V|
 &<\operatorname{perm}V
       \left\{(1+2^{-9})^2(1+2^{-20})^3-1\right\}\\
 &<\frac{\operatorname{perm}V}{240}<\det V.
\end{align*}
The two final inequalities are exact rational inequalities.
In particular $\det E>0$. The extracted row and column factors,
the product of the extracted signs, and $\det T_5$ are positive.
Equation~\eqref{gt:gamma-bridge} proves $f_{b,5}>0$.
\end{proof}

The distinction between the two error budgets is substantive: using
the worst entrywise error five times gives a weaker threshold.
The general estimate in the predecessor was $B(5)=635656$;
the preceding argument reduces the required finite interval to
$b<73728$, while using the same inherited phase remainder.

\subsection{An exact certificate for every remaining integer}

The cotangent coefficients satisfy
\begin{equation}\label{gt:beta-recurrence}
 \beta_1=\frac16,\qquad
 2(2n+1)\beta_n=\sum_{k=1}^{n-1}\beta_k\beta_{n-k}
                                      \quad(n\ge2).
\end{equation}
This follows from the Riccati equation for $(z/2)\cot(z/2)$.
We use the prime
\[
                         \ell=998244353=119\cdot2^{23}+1.
\]
The supplied verifier proves primality by trial division. It also
verifies that $3$ is a primitive root by checking the three prime
divisors $2,7,17$ of $\ell-1$. Every polynomial transform used in
the certificate has length at most $2^{18}$, so all required roots
of unity exist in $\mathbb F_\ell$.

For $6\le b\le73727$, the largest needed fundamental row is
$73732$ and the largest needed cotangent index is $73731$.
All recurrence denominators are invertible modulo $\ell$.
This is immediate for $2(2n+1)$ in \eqref{gt:beta-recurrence};
for the other diagonal use
\begin{equation}\label{gt:d-factor}
 d_s=2\frac{(2s)!}{(2s-10)!}\frac{2s-5}{s},\qquad s\ge6.
\end{equation}
Every displayed integer factor is positive and less than $\ell$.
Induction therefore identifies the computed residues with the
reductions of the exact rational coefficients.

\begin{proposition}[Finite gap-ten certificate]
\label{gt:finite}
For every integer $6\le b\le73727$, the rational determinant
$f_{b,5}$ has a nonzero reduction modulo $998244353$.
\end{proposition}

\begin{proof}[Exact computational proof]
The program \texttt{verify\_gap\_ten.cpp} constructs the necessary
cotangent coefficients and five fundamental columns using integer
arithmetic only. It then evaluates all $73722$ five-by-five
determinants by Gaussian elimination over $\mathbb F_\ell$.
Every residue is nonzero. The full indexed list is
\texttt{gap\_ten\_certificate.csv}.
\begin{samepage}
\noindent Its SHA-256 identifier is
\begin{center}\small\ttfamily
f3f749a4c4e02c27b32a58e1f65714bab568a08f3b0ebdb0fd8cb8d4bf74184d
\end{center}
\end{samepage}
The recurrence and the nonzero residues are the certificate; its
hash serves only to identify the supplied record.

Every fundamental recurrence is independently checked after the
causal computation by two full polynomial products per column.
The cotangent coefficients are independently checked against
\eqref{gt:beta-recurrence}, after their initial construction by
formal division of cosine by normalized sine.
The separate Python verifier, which does not call the NTT code,
recomputes $506$ boundary determinants with the direct quadratic
recurrence, checks $64$ cotangent coefficients against rational
Bernoulli numbers, and compares the original Bernoulli matrix
\eqref{gt:source-entry} and its full normalization
\eqref{gt:bridge} for every $6\le b\le18$.
All these checks pass; their exact values are supplied in
\texttt{gap\_ten\_independent\_checks.json}.
\end{proof}

Propositions~\ref{gt:positive-tail} and \ref{gt:finite}, followed
by \eqref{gt:bridge} and \eqref{gt:prefix-det}, prove
Theorem~\ref{gt:main}. The finite-field part proves nonvanishing,
not the real signs of the corresponding rational numbers.
Indeed direct rational checks give $f_{6,5}>0$, $f_{7,5}<0$,
and $f_{18,5}>0$. A blanket positivity claim at every admissible
index would be false.

\subsection{Why the finite recurrence can be verified efficiently}
\label{gt:fast}

The accelerated calculation has a short exact algebraic description.
Define formal series
\[
 S(z)=\sum_{n\ge0}\frac{(-1)^n z^n}{4^n(2n+1)!},\qquad
 C(z)=\sum_{n\ge0}\frac{(-1)^n z^n}{4^n(2n)!}.
\]
Since $S(0)=1$, its inverse exists over $\mathbb F_\ell$, and
\[
             1-\tfrac12\sum_{n\ge1}\beta_nz^n=C(z)/S(z).
\]
Newton iteration $U\mapsto U(2-SU)$ doubles the number of correct
inverse coefficients. The number-theoretic transform computes each
polynomial product exactly in the finite field; it introduces no
roundoff error.

For a fundamental column $r_u=R_{u,j}$, the recurrence is
\[
 d_sr_s=b_s\sum_{u<s}\beta_{s-u}r_u
                  +\sum_{u<s}\beta_{s-u}a_ur_u.
\]
Both sums are causal convolutions. Divide an index interval into
two equal halves, compute its left half, add both left-to-right
convolutions, and then compute its right half. Induction on the
interval length proves that every already determined coefficient
contributes exactly once to every later equation. The transforms
for a given interval length use the same cotangent kernel and can
be reused across all five columns.

In a cyclic transform of length $L$, the left input is supported
in $[0,L/2)$ and the kernel in $[0,L)$. All wrapped terms land in
$[0,L/2)$; none contaminates the right half being updated.
This proves the validity of the shorter cyclic convolution used
by the implementation. For a fixed half-gap $m$, the resulting
cost is $O(mN\log^2N+Nm^3)$ field operations and $O(mN)$ stored
field elements, using ordinary cubic elimination for the consecutive
boundary determinants. The same method applies to any prescribed
half-gap; it is not tied to five columns.

\section{Inverse harmonic coordinates from bounded-size determinants}
\label{gt:green-section}

For this section let $m\ge1$ be arbitrary, $b>m$, and use the same
definitions with $10$ replaced by $2m$ and $5$ by $m$.
Thus $p=2b-2m$, $q=2b$, $F_{b,m}=(R_{b+i,j})_{i,j=1}^m$, and
assume $f_{b,m}=\det F_{b,m}\ne0$. This assumption is now a theorem
for every $b>5$ when $m=5$.

For an integer forcing index $N>m$, define the causal Green column
$\mathsf G_{s,N}$ by zero initial values for $1\le s\le m$ and
\begin{equation}\label{gt:Green}
 d_s\mathsf G_{s,N}
 =\delta_{sN}+\sum_{u<s}\beta_{s-u}(b_s+a_u)\mathsf G_{u,N}
                                  \qquad(s>m).
\end{equation}
In particular $\mathsf G_{s,N}=0$ for $s<N$ and
$\mathsf G_{N,N}=1/d_N$. Put
\[
 \mathbf r_n=(R_{n,1},\ldots,R_{n,m}),\qquad
 \mathbf g_N=(\mathsf G_{b+1,N},\ldots,
                         \mathsf G_{b+m,N})^{\mathsf T}.
\]

\begin{theorem}[Every inverse coefficient in dimension $m+1$]
\label{gt:inverse}
For $0\le k,r<b$, set $n=b-k$ and $N=r+m+1$. Then
\begin{equation}\label{gt:inverse-formula}
 \boxed{\displaystyle
 (M_{b,m}^{-1})_{k,r}
 =\frac{2(-1)^{n+N}(2n)!}{(2r+1)!}\,
 \frac{\det\begin{pmatrix}
       F_{b,m}&\mathbf g_N\\
       \mathbf r_n&\mathsf G_{n,N}
                  \end{pmatrix}}{\det F_{b,m}}.}
\end{equation}
The determinant in the numerator has order $m+1$, independent of
$b$. Every entry is rational and specified by a finite recurrence.
\end{theorem}

\begin{proof}
Let $J_{m,b}$ be the coefficient matrix of the cotangent operator
\eqref{gt:null}, with $10$ replaced by $2m$, acting on
$\sum_{n=1}^b r_nz^{2n}$ and with output rows
$N=m+1,\ldots,b+m$. Its equation with unit forcing at $N$ is
the recurrence \eqref{gt:Green}, with the boundary conditions
$r_{b+1}=\cdots=r_{b+m}=0$ imposed after truncation.
The general full solution is
\[
          r_s=\mathbf r_s\mathbf c+\mathsf G_{s,N}.
\]
Its $m$ terminal conditions give
$F_{b,m}\mathbf c=-\mathbf g_N$. Hence
\[
 (J_{m,b}^{-1})_{n-1,r}
 =\mathsf G_{n,N}-\mathbf r_nF_{b,m}^{-1}\mathbf g_N
 =\frac{\det\begin{pmatrix}F_{b,m}&\mathbf g_N\\
                  \mathbf r_n&\mathsf G_{n,N}\end{pmatrix}}
        {\det F_{b,m}}.
\]

To restore the original scales, let $K_{m,b}$ be the analogous
hyperbolic-cotangent coefficient matrix. If $M^{\mathrm{rev}}$
denotes column reversal, direct comparison with
\eqref{gt:source-entry} gives
\[
 K_{m,b}=\operatorname{diag}_{0\le r<b}
           \left(\frac2{(2r+1)!}\right)
         M_{b,m}^{\mathrm{rev}}
         \operatorname{diag}_{1\le n\le b}\bigl((2n)!\bigr).
\]
Since $2B_{2h}/(2h)!=(-1)^{h-1}\beta_h$ for $h\ge1$,
its entries satisfy
$(K_{m,b})_{r,n-1}=(-1)^{N+n}(J_{m,b})_{r,n-1}$.
The same row and column signs therefore occur in the inverse.
Finally the original row $k$ corresponds to the reversed index
$n=b-k$. These factors give precisely \eqref{gt:inverse-formula}.
\end{proof}

This formula constructs each inverse row without expanding a
$b$-square determinant. It does not remove the need to compute the
finite Green column, and it makes no minimality claim about the
determinant size. The independent verifier checks all $293$ entries
of the resulting inverse matrices at $b=6,7,8,12$ directly against
the original source matrices, and separately checks an actual
six-by-six bordered determinant.

\section{Ordinary polygamma identities at every gap of ten}
\label{gt:identities}

Let $b\ge6$, $p=2b-10$, $q=2b$, and let
$\lambda_{b+r}=(M_{b,5}^{-1})_{0,r}$ for $0\le r<b$.
Define
\begin{align}
 K_b(y)&=\sum_{j=b}^{2b-1}\lambda_jy^{2j+1},\notag\\
 P_b(x)&=K_b(x+\tfrac12)-K_b(x-\tfrac12)
                   =\sum_{r=0}^{q-1}w_rx^{2r}.\label{gt:P}
\end{align}
By Theorem~\ref{gt:inverse}, all its coefficients can be given by
six-by-six determinants and the five-column recurrence.
With $c_{v,h}$ from \eqref{gt:tail-expansion}, put
\[
 u_h=\sum_{e=0}^{h}c_{p,e}c_{q,h-e}\quad(0\le h\le5),
 \qquad
 S_b(x)=\sum_{h=0}^{5}
       \left(\sum_{e=0}^{h}u_e w_{q-1-h+e}\right)x^{10-2h}.
\]

\begin{corollary}[An all-order rational family]
\label{gt:ordinary}
For every integer $b\ge6$,
\begin{equation}\label{gt:ordinary-sum}
 \boxed{\displaystyle
 \sum_{n=0}^{\infty}\left\{
 \frac{P_b(n+1/2)}{(p-1)!(q-1)!}
   \psi^{(p-1)}(n+1)\psi^{(q-1)}(n+1)-S_b(n+1/2)
                         \right\}=-\frac12.}
\end{equation}
The sum is absolutely convergent. The polynomial $P_b$ has exact
degree $2q-2$; its leading coefficient is
\begin{equation}\label{gt:leading}
 [x^{2q-2}]P_b(x)=
       \frac{4b(4b-1)f_{b-1,5}}{d_{b+5}f_{b,5}}\ne0,
\end{equation}
where at the initial index one uses the auxiliary recurrence value
\begin{equation}\label{gt:auxiliary}
 f_{5,5}=
 \frac1{363698159177574318706747080134319581355520819200000000000}.
\end{equation}
\end{corollary}

\begin{proof}
The inverse row gives $\sum_{j=b}^{2b-1}\lambda_jC_j=Z_0=-1/2$.
Expanding the centered difference in \eqref{gt:P} and using
\eqref{gt:C} gives $\sum_rw_rD_r=-1/2$.
By \eqref{gt:tail-expansion}, the nonnegative powers in
$P_b(x_n)A_p(n)A_q(n)$ are exactly $S_b(x_n)$, and the remaining
summand is $O(x_n^{-2})$. The continued sums of the subtraction
monomials vanish because $\zeta(-2h,1/2)=0$ for every $h\ge0$.
Equivalently, near spectral parameter zero, subtracting this finite
polynomial gives a normally convergent Dirichlet series and leaves
only those Hurwitz-zeta terms. Its value at zero is the ordinary
sum of the subtracted terms. Finally
$A_v(n)=\psi^{(v-1)}(n+1)/(v-1)!$ for positive even $v$.

For $b\ge7$, delete the last row and first column of $M_{b,5}$.
Reindexing the two orders down by two and the two source indices
down by one identifies the minor exactly with $M_{b-1,5}$,
including the omitted-binomial convention. Cramer's rule and
\eqref{gt:bridge} give
\[
 \lambda_{2b-1}=(-1)^{b-1}\frac{\det M_{b-1,5}}{\det M_{b,5}}
               =\frac{4b f_{b-1,5}}{d_{b+5}f_{b,5}}.
\]
The centered difference multiplies this highest coefficient by
$4b-1$, proving \eqref{gt:leading}. At $b=6$ direct rational
recurrence gives \eqref{gt:auxiliary}, and the original matrix
gives $\lambda_{11}=132/127$, verifying the same formula without
introducing a zero-order tail. Nonvanishing follows from
Theorem~\ref{gt:main} and the explicit auxiliary value.
\end{proof}

\subsection[The newly covered pair (14,24)]{The newly covered pair $(14,24)$}

Take $b=12$, so $p=14$ and $q=24$. Formula \eqref{gt:ordinary-sum}
is then the explicit ordinary identity
\[
 \sum_{n\ge0}\left\{
 \frac{P_{12}(n+1/2)}{13!\,23!}
 \psi^{(13)}(n+1)\psi^{(23)}(n+1)-S_{12}(n+1/2)
                         \right\}=-\frac12,
\]
with $\deg P_{12}=46$ and $\deg S_{12}=10$. Its leading coefficients are
\begin{align*}
 [x^{46}]P_{12}(x)
 &=\frac{196667260793027525874874688960}
          {14009065684434488629739697248281},\\
 [x^{10}]S_{12}(x)
 &=\frac{196667260793027525874874688960}
          {4188710639645912100292169477236019}.
\end{align*}
The second denominator is $13\cdot23$ times the first.
The file \texttt{gap\_ten\_14\_24\_example.json} contains all $24$
even coefficients of $P_{12}$, all six coefficients of $S_{12}$,
and the twelve coefficients of $K_{12}$ as exact rationals.
The formula \eqref{gt:inverse-formula} is an alternative finite
specification of the same coefficients. The generator checks the
inverse against the original Bernoulli matrix and the entire
subtraction polynomial symbolically. This pair lies beyond the
previous smaller-order-through-twelve coverage.

More generally, replace row zero by row $k$ of the inverse in
\eqref{gt:P}. The identical construction, with the same degree
bound and subtraction rule, gives the right-hand side
$\zeta(2k)$ for every $0\le k<b$. Thus the completed reconstruction
theorem gives a family of ordinary convergent polygamma-product
representations for every even zeta coordinate in the source block.

% ===== 04_gamma_primitives.tex =====
\section{Closed Gamma primitive certificates and analytic infinite jets}
\label{hp:section}

The newer of the two incoming \emph{Universal Gamma Normal Forms}
packages gives every canonical coefficient of the inverse-Gamma
differential hierarchy. Its research questions~1 and~6 ask, respectively,
for a nonrecursive formula for the compensating total derivative and
for analytic meaning of the infinite-jet generator. We answer those
questions here. The algebraic device used for the first answer is the
classical one-dimensional Euler homotopy operator
\cite{PooleHereman,HeremanDeconinckPoole}. Its application to the incoming
coefficient theorem, rather than the homotopy operator itself, is the
continuation result. The convergence theorem below supplies an analytic
domain for actual holomorphic logarithmic derivatives.

\subsection{Definitions and the inherited coefficient formula}

Let
\[
 \mathcal R=\mathbb Q[u_0,u_1,\ldots],\qquad Du_j=u_{j+1},
 \qquad \operatorname{wt}(u_j)=j+1.
\]
The factor degree of $u^{\boldsymbol a}=\prod_j u_j^{a_j}$ is
$|\boldsymbol a|=\sum_j a_j$. A monomial is canonical if it is a
power of $u_0$, or if its largest positive index occurs at least twice.
The canonical complement of $D\mathcal R$ gives unique polynomials
$\mathcal N(P),\mathcal H(P)$ satisfying
\begin{equation}\label{hp:decomposition}
 P=\mathcal N(P)+D\mathcal H(P),\qquad
 [1]\mathcal H(P)=0.
\end{equation}
This complement is classical \cite[Proposition~4.10]{Guo2014}.
For completeness, a unique largest derivative $u_j$ can be removed by
writing a suitable multiple of the monomial with $u_j$ replaced by
$u_{j-1}$ as a total derivative. The remainder has lower highest
index. This proves existence by termination. If $H$ is nonconstant
with highest index $j$, then $DH$ contains the nonzero part
$u_{j+1}\partial H/\partial u_j$, whose monomials are noncanonical.
Thus a canonical total derivative is zero, proving uniqueness.

Define
\begin{equation}\label{hp:hierarchy}
 A_0=1,\qquad A_{m+1}=u_0A_m-DA_m,
 \qquad N_m=\mathcal N(A_m),\quad H_m=\mathcal H(A_m).
\end{equation}
Their raw generating series is
\begin{equation}\label{hp:raw-generator}
 \sum_{m\ge0}A_m\frac{t^m}{m!}
 =\exp\left(\sum_{j\ge0}
                   \frac{(-1)^j u_jt^{j+1}}{(j+1)!}\right).
\end{equation}
Consequently, at weight $m=\sum(j+1)a_j$,
\begin{equation}\label{hp:raw-coefficient}
 a(\boldsymbol a):=[u^{\boldsymbol a}]A_m
 =\frac{(-1)^{m-|\boldsymbol a|}m!}
        {\prod_j\bigl(a_j!\,((j+1)!)^{a_j}\bigr)}.
\end{equation}

Here is the precise part of the incoming theorem used below
\cite{UniversalIncoming}. Set
\begin{equation}\label{hp:C}
 C_j(v)=\frac{(v+1)^{j+1}-(v-1)^{j+1}}
                  {2^{j+1}(j+1)!}.
\end{equation}
The coefficient $n(\boldsymbol a)=[u^{\boldsymbol a}]N_m$ is one
for $u_0^m$, zero for a unique highest positive derivative, and zero
when $\sum j a_j$ is odd. In the remaining case, with highest index
$d\ge1$ and $a_d\ge2$, it is
\begin{equation}\label{hp:normal-coefficient}
 n(\boldsymbol a)=
 \frac{(-1)^d m!}{4^d d!(d+1)!\prod_j a_j!}
 \int_0^1(1-v^2)^d
     C_d(v)^{a_d-2}\prod_{j=1}^{d-1}C_j(v)^{a_j}\,dv.
\end{equation}
The integrand is a polynomial, so this is an explicit finite rational
number. We retain attribution for this coefficient formula; the new
primitive formula starts with it as an input.

\subsection{A nonrecursive formula for every compensator}

For a differential polynomial define
\begin{align}
 E(P)&=\sum_{j\ge0}(-D)^j\frac{\partial P}{\partial u_j},
                                      \label{hp:Euler}\\
 B(P)&=\sum_{j\ge1}\sum_{k=0}^{j-1}
           u_{j-1-k}(-D)^k\frac{\partial P}{\partial u_j}.
                                      \label{hp:boundary}
\end{align}
Let $K$ denote the factor-degree operator:
$K(u^{\boldsymbol a})=|\boldsymbol a|u^{\boldsymbol a}$.
Its inverse is defined on polynomials without constant term by
division of each coefficient by its positive factor degree.

\begin{lemma}[Euler homotopy identity]\label{hp:homotopy-lemma}
For every $P\in\mathcal R$,
\begin{equation}\label{hp:Euler-boundary}
                 KP=u_0E(P)+DB(P).
\end{equation}
If $P$ has zero constant term and $E(P)=0$, then
\begin{equation}\label{hp:closed-homotopy}
                   D\bigl(BK^{-1}P\bigr)=P.
\end{equation}
\end{lemma}

\begin{proof}
For an arbitrary polynomial $Q$ and $j\ge1$, direct differentiation
of a telescoping sum gives
\[
 u_jQ-u_0(-D)^jQ
 =D\left(\sum_{k=0}^{j-1}u_{j-1-k}(-D)^kQ\right).
\]
Substitute $Q=\partial P/\partial u_j$ and sum over $j$, including
the $j=0$ term separately. Euler's identity for factor degree yields
\eqref{hp:Euler-boundary}. The operators $D$ and $B$ preserve factor
degree, while $E$ lowers it by one. Apply the identity separately
to each positive factor-degree component of $P$, then divide by that
degree. This proves \eqref{hp:closed-homotopy}. In particular no
analytic integration and no recursive integration by parts is involved.
\end{proof}

\begin{theorem}[Complete finite primitive formula]\label{hp:closed}
For every $m\ge1$,
\begin{equation}\label{hp:Hclosed}
 \boxed{\displaystyle H_m=BK^{-1}(A_m-N_m).}
\end{equation}
Equivalently, let $p(\boldsymbol a)=a(\boldsymbol a)-n(\boldsymbol a)$
with the explicit coefficients in
\eqref{hp:raw-coefficient}--\eqref{hp:normal-coefficient}. Then
\begin{equation}\label{hp:finite-sum}
 \boxed{\displaystyle
 H_m=\sum_{\sum(j+1)a_j=m}\frac{p(\boldsymbol a)}{|\boldsymbol a|}
       \sum_{j\ge1}a_j\sum_{k=0}^{j-1}(-1)^k
                  u_{j-1-k}D^k u^{\boldsymbol a-\boldsymbol e_j}.}
\end{equation}
All sums are finite. Together with \eqref{hp:normal-coefficient},
this determines both polynomials in the canonical decomposition
without any recursive normal reduction.
\end{theorem}

\begin{proof}
The polynomial $A_m-N_m$ is a total derivative by
\eqref{hp:decomposition}. The identity $E(DQ)=0$ follows by cancellation
from
$\partial_{u_j}D=D\partial_{u_j}+\partial_{u_{j-1}}$.
Lemma~\ref{hp:homotopy-lemma} therefore shows that the right side of
\eqref{hp:Hclosed} differentiates to $A_m-N_m$.
It has zero constant term, so uniqueness in
\eqref{hp:decomposition} identifies it with $H_m$.
Expanding the two operators gives \eqref{hp:finite-sum}.
\end{proof}

If a formula without powers of $D$ is desired, every term is expanded by
\begin{equation}\label{hp:Dexpansion}
 D^k u^{\boldsymbol b}
 =k!\sum_{\substack{\sum_{r\ge0}b_{i,r}=b_i\ (i\ge0)\\
                         \sum_{i,r}r b_{i,r}=k}}
       \prod_i\frac{b_i!}{\prod_r b_{i,r}!(r!)^{b_{i,r}}}
       \prod_{i,r}u_{i+r}^{b_{i,r}}.
\end{equation}
This is simply the multinomial Leibniz rule. Thus the word
``nonrecursive'' describes a literal finite coefficient formula,
not an instruction to run an integration-by-parts recurrence in disguise.

\subsection{An explicit higher Gamma primitive}

For example,
\begin{align}
 N_6={}&u_0^6-15u_0^2u_1^2+u_2^2,\label{hp:N6}\\
 H_6={}&-3u_0^5+20u_0^3u_1-15u_0^2u_2-15u_0u_1^2
             +6u_0u_3+9u_1u_2-u_4.\label{hp:H6}
\end{align}
At weight eight the normal part is
\begin{equation}\label{hp:N8}
 N_8=u_0^8-70u_0^4u_1^2+28u_0^2u_2^2-7u_1^4-u_3^2,
\end{equation}
and the companion polynomial, supplied completely by the same formula, is
\begin{align}
 H_8={}&-4u_0^7+56u_0^5u_1-70u_0^4u_2-140u_0^3u_1^2
              +56u_0^3u_3\notag\\
      &+252u_0^2u_1u_2-28u_0^2u_4+112u_0u_1^3
              -112u_0u_1u_3\notag\\
      &-84u_0u_2^2+8u_0u_5-98u_1^2u_2+20u_1u_4
              +36u_2u_3-u_6.\label{hp:H8}
\end{align}
These low-weight examples are illustrations of the closed construction,
not claims of new historical evaluations.

For a smooth function $U$, substitute $u_j=U^{(j)}$ and let $Y'=-UY$.
The inherited weighted identity now has a fully nonrecursive right side:
\begin{equation}\label{hp:weighted}
 \boxed{\displaystyle
 \int Y\bigl(N_m(U)+U H_m(U)\bigr)\,dx
       =-Y\bigl(A_{m-1}(U)+H_m(U)\bigr)+C.}
\end{equation}
Indeed $(YA_{m-1})'=-YA_m$ and
$(YH_m)'=Y(DH_m-UH_m)$. The formula follows on adding these two
identities. The compensator $UH_m$ and the endpoint value of $YH_m$
are part of the formula; neither can be discarded.

Taking $U=\psi$ and $Y=1/\Gamma$ gives exact Gamma--polygamma
antiderivatives at every weight. Taking instead
\begin{equation}\label{hp:StieltjesU}
 F_q(x)=\frac{(-1)^q}{q+1}\zeta^{(q+1)}(0,x),\quad
 U_q(x)=-\gamma_q(x),\quad Y_q(x)=e^{-F_q(x)}
\end{equation}
gives the corresponding all-index Stieltjes primitives.
Here $\zeta^{(q+1)}$ denotes an order derivative, not an argument
derivative. The standard identity
$\partial_x\zeta(s,x)=-s\zeta(s+1,x)$ and the Laurent convention
$\zeta(1+s,x)=s^{-1}+\sum_{q\ge0}(-1)^q\gamma_q(x)s^q/q!$
give $F_q'=-\gamma_q=U_q$ \cite{DLMFHurwitz}.
Thus $Y_q'=-U_qY_q$, as required. At $q=0$,
$F_0=\log\Gamma-\tfrac12\log(2\pi)$.

\subsection{The infinite-jet formula is analytic for holomorphic fields}

The incoming coefficient theorem is formal in infinitely many jets.
The following result permits their actual substitution from a
holomorphic function. Define the entire function
\[
 \phi(z)=\frac{e^z-1-z}{z^2},\qquad \phi(0)=\frac12.
\]
For $x$ in a domain of holomorphy of $U$, put
$x_j=t^{j+1}U^{(j)}(x)$ and
$S_{r-1}(v)=\sum_{j=1}^{r-1}C_j(v)x_j$.

\begin{theorem}[Analytic normal generator]\label{hp:analyticN}
Suppose $U$ is holomorphic on $|x-x_0|<R$. The series
\begin{equation}\label{hp:Nanalytic}
 \boxed{\displaystyle
 \mathfrak N_U(x,t)=e^{tU(x)}\left[
 1+\frac12\sum_{r\ge1}
 \frac{(-1)^r x_r^2}{4^r r!(r+1)!}
 \int_{-1}^1(1-v^2)^r e^{S_{r-1}(v)}
                         \phi(C_r(v)x_r)\,dv\right]}
\end{equation}
converges absolutely and locally uniformly when
\begin{equation}\label{hp:normaldomain}
                        |x-x_0|+|t|<R.
\end{equation}
It is holomorphic there and has Taylor series
\begin{equation}\label{hp:normalTaylor}
            \mathfrak N_U(x,t)=\sum_{m\ge0}N_m(U(x))\frac{t^m}{m!}.
\end{equation}
\end{theorem}

\begin{proof}
Cover any compact subset of \eqref{hp:normaldomain} by finitely
many sufficiently small compact polydisks. On each such polydisk
choose $\rho$ with $\sup|t|<\rho<R-\sup|x-x_0|$. Cauchy's
estimate supplies $M>0$ such that, uniformly there,
\[
 |U^{(r)}(x)|\le Mr!\rho^{-r},\qquad |t|<\rho.
\]
For real $-1\le v\le1$, the explicit positive-coefficient formula
for $C_j$ gives $|C_j(v)|\le1/(j+1)!$. Hence
\[
 \sum_{j\ge1}|C_j(v)x_j|
 \le M|t|\sum_{j\ge1}\frac{(|t|/\rho)^j}{j+1}<\infty
\]
uniformly. Also $|\phi(z)|\le e^{|z|}/2$, from its elementary
integral representation
$\phi(z)=\int_0^1(1-s)e^{sz}\,ds$.
Apart from a uniform constant, the absolute value of the $r$th
summand of \eqref{hp:Nanalytic} is consequently at most
\[
 \frac{|x_r|^2}{4^r r!(r+1)!}
 \le\frac{M^2|t|^2}{r+1}
                  \left(\frac{|t|}{2\rho}\right)^{2r}.
\]
This is summable. The same reasoning applies on each compact
subset, proving normal convergence and holomorphy. At any fixed
power of $t$ only finitely many $r$ contribute. Expanding the
exponential and $\phi$, then using
\eqref{hp:normal-coefficient}, gives exactly the canonical
coefficients. This proves \eqref{hp:normalTaylor}.
\end{proof}

The factor $\phi$ removes all apparent singularities before any
infinite summation. In particular odd-index kernels cause no
finite-part integral in the analytic formula.

\begin{theorem}[Analytic primitive and translation identity]
\label{hp:analyticH}
Under the same hypothesis, the canonical series
\begin{equation}\label{hp:Hanalytic}
           \mathfrak H_U(x,t)=\sum_{m\ge1}H_m(U(x))\frac{t^m}{m!}
\end{equation}
converges locally uniformly on the smaller domain
\begin{equation}\label{hp:primitivedomain}
                        |x-x_0|+2|t|<R.
\end{equation}
If $F'=U$ locally, then on this domain
\begin{equation}\label{hp:translation}
 \boxed{\displaystyle
 e^{F(x)-F(x-t)}=\mathfrak N_U(x,t)
                               +\partial_x\mathfrak H_U(x,t).}
\end{equation}
The expression on the left is independent of the additive constant
in $F$.
\end{theorem}

\begin{proof}
We provide the convergence justification, since the formal
decomposition alone is insufficient. Let $P=A-N$ denote the
formal generating-series difference. The coefficientwise homotopy
formula can be written as
\begin{equation}\label{hp:Hintegral}
 \mathfrak H_U=\int_0^1\sum_{j\ge1}\sum_{k=0}^{j-1}
 U^{(j-1-k)}(x)(-\partial_x)^k
       \left.(\partial_{u_j}P)\right|_{u_r=\lambda U^{(r)}(x)}
                                      \,d\lambda.
\end{equation}
The factor of $\lambda$ from differentiating $P(\lambda u)$
cancels the $1/\lambda$ in $K^{-1}P=\int_0^1P(\lambda u)d\lambda/\lambda$.
There is therefore no singular endpoint in \eqref{hp:Hintegral}.

Fix a compact subset of \eqref{hp:primitivedomain}. Around each
of its $x$ values choose a disk of radius $\rho>|t|$ whose
closure still satisfies the normal-generator condition with
the same $t$. The raw exponential gives immediately
\[
 |\partial_{u_j}A|\le C\frac{|t|^{j+1}}{(j+1)!}.
\]
Differentiation of \eqref{hp:Nanalytic}, uniformly for
$0\le\lambda\le1$, gives the same bound for
$\partial_{u_j}N$. To see the only potentially different term,
differentiate the $r=j$ summand in its nonsingular form and use
\[
 \left|\frac{d}{dy}\{y^2\phi(cy)\}\right|
   =\left|y\frac{e^{cy}-1}{cy}\right|
   \le |y|e^{|cy|},
\]
with continuous interpretation at $cy=0$. The prefactor
$1/(4^jj!(j+1)!)$ and Cauchy's jet estimate then give
the required bound. For $r>j$ there is just the factor
$t^{j+1}C_j$, and the already summable majorant applies.
These bounds hold on a slightly larger compact set, so Cauchy's
estimate in $x$ yields
\[
 \left|\partial_x^k(\partial_{u_j}P)(\lambda U)\right|
 \le C\frac{k!\,|t|^{j+1}}{(j+1)!\rho^k},\qquad
 |U^{(j-1-k)}(x)|\le M\frac{(j-1-k)!}{\rho^{j-1-k}}.
\]
If necessary shrink $\rho$ slightly to have both bounds with
one common radius still exceeding $|t|$. Since
\[
 \sum_{k=0}^{j-1}\frac{k!(j-1-k)!}{(j+1)!}\le\frac1{j+1},
\]
the $j$th group in \eqref{hp:Hintegral} is bounded by
$CM|t|^{j+1}\rho^{1-j}/(j+1)$, a summable series.
The same bounds prove normal convergence on compact subsets.
Coefficient extraction is justified and reproduces
Theorem~\ref{hp:closed}. Finally the convergent raw series is
$\exp(\int_{x-t}^xU(s)ds)=e^{F(x)-F(x-t)}$.
The coefficientwise equality $A_m=N_m+DH_m$, now between
holomorphic functions, proves \eqref{hp:translation}.
\end{proof}

\begin{corollary}[Gamma and Stieltjes translation formulas]
\label{hp:specialtranslations}
For positive $x$ and $|t|<x/2$,
\begin{align}
 \frac{\Gamma(x)}{\Gamma(x-t)}
     &=\mathfrak N_{\psi}(x,t)
                         +\partial_x\mathfrak H_{\psi}(x,t),
                                                    \label{hp:Gammatrans}\\
 \exp\bigl(F_q(x)-F_q(x-t)\bigr)
     &=\mathfrak N_{-\gamma_q}(x,t)
                         +\partial_x\mathfrak H_{-\gamma_q}(x,t),
                                                    \label{hp:Stieltjestrans}
\end{align}
with $F_q$ as in \eqref{hp:StieltjesU}. These are ordinary
holomorphic identities in $t$ and local holomorphic identities in $x$.
On a positive interval $[a,b]$ with $|t|<a/2$ they may be
integrated termwise and retain the exact endpoint difference
$\mathfrak H_U(b,t)-\mathfrak H_U(a,t)$.
\end{corollary}

\begin{proof}
Both $\psi$ and $\gamma_q$ are holomorphic on the right half-plane.
At a fixed positive $x$, every disk of radius $R<x$ is available.
Take $2|t|<R<x$ in Theorem~\ref{hp:analyticH}, then substitute
the standard primitives. Compactness gives uniformity on $[a,b]$.
\end{proof}

No sharp convergence radius is asserted. The purpose of the domain is
to justify infinite-jet substitution, coefficientwise differentiation,
and anchored integration for actual special functions.

\subsection{Independent finite checks and the remaining question}

The accompanying \texttt{normal\_forms.py} constructs $A_m$ by
its differential recurrence, constructs $N_m$ independently from
the rational beta coefficient formula, and computes $H_m$ in two
ways: the direct canonical reduction and the closed Euler homotopy.
Through weight~20 all results agree, including $A_m=N_m+DH_m$
and the full Euler--boundary identity. There are 106 recorded
exact assertions, including a rejected coefficient corruption.
These checks support implementation correctness; the proofs of
the all-order assertions are the theorems above.

An additional exploratory calculation through weight~24 suggests
that for $2\le m\le24$ every monomial of weight $m-1$ occurs in
$H_m$, with sign
$(-1)^{m-|\boldsymbol a|}$ for its coefficient. We do not use or
claim this as a theorem. A positivity formula for the absolute
coefficients in \eqref{hp:finite-sum} would strengthen the
closed primitive result and explain the observed full support.

% ===== 05_audit_and_questions.tex =====
\section{Source audit and proposed integration corrections}
\label{audit:section}

The audit found no new counterexample to a theorem in the inherited
material on which the proofs here depend. It did find several places
where the status of a direction must be updated when the newer reports
are integrated. An older question remains a valid historical record
even when a later incoming article answers it. The appropriate change
is a dated completion note and a link to the proof, while retaining
the original provenance.

\paragraph{A verified editorial duplication.}
The retrieved manuscript \file{README.md} repeats the paragraph beginning
``The complete 1,032-input first-letter X family consecutively.
Remove its second occurrence while retaining the first and the surrounding
research-status text. This duplication was already noted in the latest
incoming source-status report; it is an editorial cleanup, not a new
mathematical erratum. The separate statements about pending canonical
analytic integration remain appropriate and should not be rewritten as
though every incoming result had already been incorporated.

\subsection{What should change in the current research map}

The fixed-gap map should now record complete gaps $2,4,6,8,10$.
The gap-eight proof belongs to \cite{PolygammaResonances}; only gap ten
is completed in this article. The smaller-order-through-twelve results
from earlier work are another source of coverage and should not be
counted again as new gap-ten cases. The pair $(14,24)$ is included as
an explicit example outside that earlier smaller-order range.

The compact inverse-row question receives a constructive bounded-size
answer in Theorem~\ref{gt:inverse}: at fixed half-gap $m$, every entry
is a quotient involving an $(m+1)$-square determinant. This does not
answer the stronger request for a short product or for a recurrence
using only a fixed number of final rows. Nor does the gap-ten proof
remove the finite computational component; the source's request for
a structural proof without finite completion remains open.

The three-cluster question $B^r,C,D$ in
\emph{Gamma Normalizations and Colored Moments} is answered by the
explicit dictionary in Section~\ref{tri:sec:main}
\cite{GammaColored}. A status note should retain the distinction between
the finite colored dictionary for one spectral term and the convergent
infinite spectral sum for the Gamma face. The much broader cluster
pattern $B^r,C^s,D^t$ with arbitrary multiplicities remains open in the
form posed there.

The collision question should now distinguish three levels. The
two-factor theorem for all positive monomial scales and all Stieltjes
indices is inherited. The present arbitrary-factor theorem evaluates
the complete local correction at Stieltjes index zero, uniformly for
all positive paths. Complete scalar values are additionally obtained
for the one-derivative centered family, including the uncentered
triple. The corresponding arbitrary-factor theorem at positive
Stieltjes indices is a further question.

Finally, the first and sixth research questions of the expanded
\emph{Universal Gamma Normal Forms} package receive the answers in
Section~\ref{hp:section}. The universal normal coefficients themselves
remain an attributed input. The new finite formula is for the
compensating polynomial $H_m$, and the new analytic result is the
convergence of the infinite substitutions for a holomorphic field.

\subsection{Conventions that materially affect correctness}

\paragraph{Nonvanishing does not imply positivity.}
The real fundamental determinant has
$f_{6,5}>0$, $f_{b,5}<0$ for $7\le b\le17$, and $f_{18,5}>0$.
Thus the conjectural strengthening $f_{b,5}>0$ at every admissible
index is false. This is an explicit obstruction to that possible
heuristic, not a claim that the predecessor asserted it. The
finite-field certificate proves nonvanishing. Real positivity is
proved separately only for $b\ge73728$.

\paragraph{The constant depends on the scale relation.}
The nonresonant triple in \eqref{intro:triple} cannot be applied to
$(\alpha_1,\alpha_2,\alpha_3)=(3,2,1)$.
The extra terms in \eqref{pc:integer-triple} are a concrete correction
to such an extension. A power--logarithm constant is also different
from the ordinary limit of the original divergent family. The
explicit $Q$ subtraction states exactly which ordinary limit is
available.

\paragraph{Exact regular coefficients are necessary for arbitrary paths.}
Replacing $A_{S,d}(\boldsymbol a)$ by a prematurely truncated expression
in the shifts can destroy uniformity. For example, if
$a_{\min}=e^{-1/t^2}$ and $a_{\max}=t$, then
$a_{\max}|\log a_{\min}|=1/t$ does not tend to zero.
Theorem~\ref{pc:uniform} keeps the shift dependence of every retained
regular coefficient exactly; its remainder proof therefore needs no
assumption on relative rates.

\paragraph{Labelled derivatives and confluent derivatives differ.}
The mixed operator $\partial_{C_1}\cdots\partial_{C_m}$ is applied
before $C_1=\cdots=C_m$ is imposed. A total derivative along that
diagonal has extra multiplicities and other partial derivatives.
The factorial $(m-1)!$ in the confluent divided difference is part
of the theorem. Likewise, the three-cluster primitive must retain
its two edge terms; its mixed derivative alone does not determine
the primitive without boundary data.

\paragraph{Endpoints and branches stay in the formulas.}
The compensator $UH_m$ and the endpoint term $YH_m$ in
\eqref{hp:weighted} are essential. Discarding a total derivative
under an unweighted finite-part integral requires its own endpoint
calculation. For the colored formulas the displayed real branch is
continued using letters outside the integration path. Iterated-integral
shuffle keeps those letters; a stuffle conversion outside the initial
disk needs additional branch bookkeeping if it introduces new
cumulative colors.

\paragraph{The arithmetic questions remain arithmetic questions.}
An invertible rational coefficient matrix does not prove independence
of its evaluated constants. A depth-three representation does not
prove that depth three is necessary. The Gaussian $S_6$ identity and
the current revised $S_8$ period candidate remain unresolved in the
inspected research programme; the present results provide no proof of
those equalities. These restrictions should remain visible in any
combined manuscript overview.

\section{Further research questions}
\label{questions:section}

The following questions are chosen for their potential to produce
identities, structural formulas and exact evaluations. The thresholds
used in the gap-ten proof are tools for nonvanishing, rather than
the main objects of interest.

\begin{enumerate}[label=\textbf{Q\arabic*.},leftmargin=*]
\item \textbf{Arbitrary unequal even harmonic tails.}
Prove or disprove $\det M_{b,m}\ne0$ for every $b>m\ge1$.
Can a transform, a sign-regular kernel, a Pad\'e interpretation or a
spectral argument account for the nonvanishing despite the finite
sign changes? The full two-parameter problem is stronger than
settling each of finitely many gaps.

\item \textbf{The next fixed gap and a reusable exact completion theorem.}
Complete gap twelve using the effective Gamma basis and the causal
convolution method. A useful further result would be an explicit
algorithm that takes $m$, derives a valid analytic threshold and
produces a fully replayable finite certificate, including a principled
choice of primes when one reduction vanishes. A zero residue is not
evidence of a zero rational determinant.

\item \textbf{Closed Green columns and inverse harmonic coordinates.}
Find a hypergeometric, Pad\'e or integrable representation of
$\mathsf G_{s,N}$ that eliminates its long causal recurrence.
Determine whether the bounded determinant in
\eqref{gt:inverse-formula} admits a short quotient of special functions.
The resulting formulas would turn high-order zeta-coordinate sums
into more recognizable identities.

\item \textbf{Positive-index Stieltjes collisions with three or more factors.}
For $q_i\ge0$, extend the uniform local theorem to
$\prod_i\partial_x^{r_i}\gamma_{q_i}(x+a_i)$, with a fully evaluated
correction and a uniform error for arbitrary positive paths.
The singular kernels now contain logarithms of several shifted
variables. An effective route is to introduce independent spectral
powers before differentiating, while retaining every regular
coefficient needed for uniformity. A formal derivative of an
unproved uniform family is insufficient.

\item \textbf{A finite classification of resonant scale contributions.}
For a fixed derivative pattern and fixed number of factors, classify
the integer relations among the $\alpha_i$ that can contribute to
the log-free constant. The rational kernel makes each specified
monomial case finite. Is there a compact chamber description with
explicit transition polynomials and amplitude logarithms?
The difference between \eqref{intro:triple} and
\eqref{pc:integer-triple} provides a first exact test case.

\item \textbf{More than one derivative-bearing factor.}
Find full scalar evaluations for centered products with two or more
nonzero derivative orders. The local correction is already given by
Theorem~\ref{pc:uniform}; the unresolved step is the merged product.
Which patterns are total derivatives of finite polygamma polynomials,
and which force marked-zeta or genuine new moment coordinates?

\item \textbf{Generating functions for the ordered constant rows.}
Use \eqref{pc:all-powers-formula} to sum in $N$ and in the derivative
position $r$. Seek polylogarithmic or Gamma expressions for these
double generating functions, and relations with harmonic-number
weights. The adjacent-difference and zero-sum identities give two
exact constraints on any proposed answer.

\item \textbf{A finite primitive for more displaced masses.}
Integrate \eqref{tri:eq:Gamma-mixed} for
$B^r,C_1,\ldots,C_m$ with $m\ge3$, keeping all lower-dimensional
face data. Is depth at most $m+1$ sufficient uniformly in $r$?
Can the edge reconstruction be expressed by a finite symmetric
dictionary with controlled letters, rather than by unresolved
parameter derivatives?

\item \textbf{Three clusters with arbitrary multiplicities.}
Extend $B^r,C,D$ to $B^r,C^s,D^t$ and identify the smallest useful
uniform depth bound. The one-unit multiplicities in the present
proof produce an elementary simplex resolvent; the general problem
requires retaining additional Dirichlet powers without losing the
finite colored description.

\item \textbf{Arithmetic collapse of the Gamma spectral sum.}
Find special positive $a$ and nonnegative masses for which
\eqref{tri:eq:Gamma} reduces to finitely many Gamma, zeta or
multiple-polylogarithm constants. Coincident masses, rational
shifts and imaginary-Gamma multiplication identities give natural
families to test. Any reduction should state its period basis and
branch conventions before invoking a numerical relation search.

\item \textbf{Signs and full support of the compensators.}
The exact data through weight $24$ suggest that for every $m\ge2$
all monomials of weight $m-1$ occur in $H_m$, with coefficient sign
$(-1)^{m-|\boldsymbol a|}$.
Prove or refute this statement. A positive integral formula for the
absolute coefficients would explain the cancellations in the
Euler-homotopy expression and could also control their denominators.

\item \textbf{Sharp analytic domains and continuation of the canonical series.}
Determine the true convergence domains of $\mathfrak N_U$ and
$\mathfrak H_U$ for $U=\psi$ and $U=-\gamma_q$.
The sufficient domain $|t|<x/2$ makes the translation identity
rigorous; it is not asserted to be sharp. Can the normal and
primitive pieces continue separately to the natural singular locus
of the Gamma or Stieltjes ratio, or do they develop additional
singularities that cancel only in their sum?

\item \textbf{Label-preserving and difference-algebra normal forms.}
Develop a canonical complement for several independent fields that
retains the individual Gamma factors and their mixed derivative
relations. In a separate direction, replace the ordinary derivation
by a difference operator and derive the changed homotopy and product
rules. Simply substituting a $q$-Gamma function into an ordinary
derivative formula would not resolve this structural question.

\item \textbf{The unresolved Gaussian periods.}
Continue the Gaussian $S_6$ and revised $S_8$ problems in a declared
relation system. Seek a new analytic transformation, a complete
iterated-integral reduction or a certified functional identity.
Finite relation searches and numerical agreement can guide that
work but do not become proofs merely by increasing precision or
the size of a relation matrix.
\end{enumerate}

% ===== 06_reproducibility.tex =====
\section{Proof certificates and reproducibility}
\label{repro:section}

The archive contains the standalone source and PDF of this article,
the same text in modular sections for integration, exact coefficient
data, numerical receipts, and the programs used to generate them.
The build is local and requires no network access. The scripts never
need repository credentials or external services.

\subsection{What was checked}

\begin{longtable}{@{}>{\raggedright\arraybackslash}p{0.22\textwidth}>{\raggedright\arraybackslash}p{0.48\textwidth}>{\raggedright\arraybackslash}p{0.23\textwidth}@{}}
\toprule
Component & Recorded check & Mathematical role\\
\midrule
\endhead
Gap-ten finite interval & $73722$ nonzero residues for every
$6\le b\le73727$ modulo $998244353$. & Exact finite part of the
proof.\\[5pt]
Gap-ten independent audit & All full recurrence identities;
$506$ direct boundary determinants; $64$ Bernoulli coefficients;
original matrices for $6\le b\le18$; $293$ inverse entries at
$b=6,7,8,12$. & Independent normalization and implementation
checks.\\[5pt]
Gap-ten infinite tail & A complex Stirling remainder and a rational
determinant/permanent margin for $b\ge73728$. & Analytic proof,
not a numerical extrapolation.\\[5pt]
Polygamma contacts & $21$ exact kernel checks, including repeated
nodes; the complete resonant triple; centered rows $N=2,\ldots,8$,
their adjacent differences and zero sums. & Exact finite checks
of algebra and implementation.\\[5pt]
Shifted-integral diagnostics & $12$ quadratures at $65$ decimal
working precision, with all three derivative positions and both
centered and uncentered triples. & Asymptotic diagnostics; not
interval certificates.\\[5pt]
Three-cluster formulas & $192$ exact harmonic checks, $14$ scalar
comparisons, $24$ mixed-derivative comparisons and one Gamma
spectral comparison: $231$ assertions in all. & Exact arithmetic
plus floating-point diagnostics, separated in the receipt.\\[5pt]
Gamma normal forms & $106$ exact assertions through weight $20$,
using an independent reducer, the beta coefficient formula and
Euler homotopy, including a rejected coefficient corruption.
& Finite verification of the all-order algebraic formulas.\\
\bottomrule
\end{longtable}

The finite-field proof is reproducible from its recurrence and
source code. Each transform is arithmetic in a prime field, with
the prime, primitive root and nonzero denominators checked.
No floating-point operation decides the nonvanishing statement.
The certificate hash in Section~\ref{gt:section} identifies the
record, while the verifier establishes the mathematical content.

The three-cluster numerical comparisons use independently arranged
formulas: cumulative-color sums, an ordinary scalar logarithmic
integral and coefficients from the Dirichlet moment expansion.
At $75$ decimal working precision the scalar scaled residuals are
below $1.5\times10^{-75}$ and the mixed-derivative residuals below
$6\times10^{-72}$ for the recorded cases. These statements describe
observed residuals, not certified errors in every reported digit.

For the collision quadratures, the largest residual after subtracting
the exact contact and the merged baseline decreases from approximately
$0.0187914$ at $\epsilon=10^{-3}$ to $0.000372220$ at
$\epsilon=10^{-5}$. The proof of the limit is uniform dominated
convergence. The quadratures are useful precisely because they compare
the original positive-shift integral to the independently assembled
finite correction.

\subsection{Programs and exact records}

The main reproduction command is
\begin{quote}
\verb|python3 reproduce.py --exact|
\end{quote}
It rebuilds the C++ verifier, recomputes the full gap-ten interval,
compares its bytes with the included certificate, and runs the
independent rational, contact and normal-form checks.
The optional numerical suite is
\begin{quote}
\verb|python3 reproduce.py --numeric|
\end{quote}
The PDF is built from the standalone source with
\begin{quote}
\verb|latexmk -pdf Higher_Identities.tex|
\end{quote}
The file \file{README.txt} gives the dependencies and individual
commands. Exact rational data use numerator/denominator strings;
they do not require a computer algebra system to parse.
The full $(14,24)$ polynomial pair is in
\file{code/gap_ten/gap_ten_14_24_example.json}.

\subsection{Provenance and limits of the verification}

The supplied source manifest records the bytes retrieved from the
moving main branch. The repository HEAD reported during retrieval was
\begin{center}\small\ttfamily
e4ea0a4c7937d7ac0091048065222d1b58370482.
\end{center}
The individual file hashes, rather than an assumption that every
request was an atomic checkout of that HEAD, are the definitive
identifiers of the inspected snapshots. The two similarly named
\emph{Universal Gamma Normal Forms} archives are both recorded;
the expanded version with the suffix \texttt{(2)} supplies the
two explicit research questions addressed here.

The principal analytic dependencies have been restated with their
normalizations. The complementary-minor and Gamma-basis constructions
in the gap-ten section retain their attribution to the ordered-boundary
report; the universal normal coefficients retain theirs to the
universal-normal-form report. Classical sources support the standard
Gamma, Hurwitz-zeta, divided-difference and homotopy facts. A separate
\file{INTEGRATION_NOTES.txt} maps the new conclusions to the existing
questions without rewriting historical archives.

The proofs have been checked across the independent research sections,
and the displayed finite formulas have the recorded computational
support. This is a research contribution for further expert review
and integration, not a claim that the entire surrounding manuscript
or the worldwide literature has been exhaustively verified.

% ===== 07_bibliography.tex =====
\begin{thebibliography}{99}
\small
\setlength{\itemsep}{3pt}
\setlength{\parsep}{0pt}
\setlength{\parskip}{0pt}

\bibitem{RepoManuscript}
V.~Reshetnikov and the ProveIt research collection,
\emph{Analysis/Polylogarithms manuscript}, retrieved 12 October 2026.
\url{https://github.com/VladimirReshetnikov/ProveIt/tree/main/Analysis/Polylogarithms/docs/manuscript}.

\bibitem{RepoIncoming}
ProveIt research collection, \emph{Incoming research archives},
retrieved 12 October 2026. The supplied source manifest identifies
all twelve retrieved archives.
\url{https://github.com/VladimirReshetnikov/ProveIt/tree/main/docs/incoming}.

\bibitem{OrderedBoundary}
ProveIt research collection,
\emph{Ordered Boundary Identities}, incoming package dated
11 October 2026, especially
\file{sections/08_gamma_basis.tex}.
Archive: \file{ProveIt_Ordered_Boundary_Identities_2026-10-11.zip},
in \cite{RepoIncoming}.

\bibitem{PolygammaResonances}
ProveIt research collection,
\emph{Polygamma Moments and Stieltjes Resonances}, incoming package
dated 12 October 2026, especially
\file{sections/02_gap_reconstruction.tex},
\file{sections/03_multiscale_collisions.tex}, and
\file{sections/05_source_status.tex}.
Archive: \file{ProveIt_Polygamma_Moments_and_Stieltjes_Resonances_2026-10-12.zip},
in \cite{RepoIncoming}.

\bibitem{GammaColored}
ProveIt research collection,
\emph{Gamma Normalizations and Colored Moments}, incoming package
dated 12 October 2026. The two-cluster dictionary and its further
research questions are the starting point for
Section~\ref{tri:sec:main}.
Archive: \file{ProveIt_Gamma_Normalizations_and_Colored_Moments_2026-10-12.zip},
in \cite{RepoIncoming}.

\bibitem{UniversalIncoming}
ProveIt research collection,
\emph{Universal Gamma Normal Forms}, expanded incoming package dated
11 October 2026, especially the universal coefficient theorem and
\file{sections/11_audit.tex}, questions 1 and 6.
Archive: \file{ProveIt_Universal_Gamma_Normal_Forms_2026-10-11 (2).zip},
in \cite{RepoIncoming}.

\bibitem{BilinearTails}
ProveIt research collection,
\emph{Bilinear Tails and Harmonic Determinants}, report in the
Stieltjes-harmonic-resonance continuation, especially the compact
inverse-row question. Source location is recorded in
\file{INTEGRATION_NOTES.txt}; see also the later reformulation in
\cite{PolygammaResonances}.

\bibitem{DLMF}
NIST Digital Library of Mathematical Functions,
Chapter 5, especially Sections 5.7, 5.8, 5.11 and 5.15.
\url{https://dlmf.nist.gov/5}.

\bibitem{DLMFHurwitz}
NIST Digital Library of Mathematical Functions,
Section 25.11, \emph{Hurwitz Zeta Function}.
\url{https://dlmf.nist.gov/25.11}.

\bibitem{DLMFPolylog}
NIST Digital Library of Mathematical Functions,
Section 25.12, \emph{Polylogarithms}.
\url{https://dlmf.nist.gov/25.12}.

\bibitem{DeBoor}
C.~de Boor, \emph{Divided Differences},
Surveys in Approximation Theory \textbf{1} (2005), 46--69.
\url{https://arxiv.org/abs/math/0502036}.

\bibitem{Lesch}
M.~Lesch, \emph{Divided Differences in Noncommutative Geometry:
Rearrangement Lemma, Functional Calculus and Expansional Formula},
Journal of Noncommutative Geometry \textbf{11} (2017), 193--223.
\url{https://arxiv.org/abs/1405.0863}.

\bibitem{TicaGalapon}
C.~D.~Tica and E.~A.~Galapon,
\emph{Finite-Part Integration of the Generalized Stieltjes Transform
and its dominant asymptotic behavior for small values of the parameter
Part I: Integer orders}, Journal of Mathematical Physics
\textbf{59}, 023509 (2018).
\url{https://arxiv.org/abs/1703.07979}.

\bibitem{VillanuevaGalapon}
L.~Villanueva and E.~A.~Galapon,
\emph{Finite-Part Integration in the Presence of Competing Singularities:
Transformation Equations for the hypergeometric functions arising from
Finite-Part Integration}, Journal of Mathematical Physics
\textbf{62}, 043505 (2021).
\url{https://arxiv.org/abs/2012.01943}.

\bibitem{PooleHereman}
D.~Poole and W.~Hereman,
\emph{The homotopy operator method for symbolic integration by parts
and inversion of divergences with applications},
Applicable Analysis \textbf{89} (2010), no.~4.
\url{https://arxiv.org/abs/0908.0399}.
\url{https://doi.org/10.1080/00036810903208155}.

\bibitem{HeremanDeconinckPoole}
W.~Hereman, B.~Deconinck and L.~D.~Poole,
\emph{Continuous and Discrete Homotopy Operators:
A Theoretical Approach made Concrete},
Mathematics and Computers in Simulation \textbf{74} (2007), 352--360.
\url{https://arxiv.org/abs/nlin/0510064}.

\bibitem{Guo2014}
L.~Guo, G.~Regensburger and M.~Rosenkranz,
\emph{On Integro-Differential Algebras},
Journal of Pure and Applied Algebra \textbf{218} (2014), no.~3,
especially Proposition~4.10.
\url{https://arxiv.org/abs/1212.0266}.
\url{https://doi.org/10.1016/j.jpaa.2013.06.015}.

\end{thebibliography}
\end{document}

